\documentclass[11pt]{article}

% ACL style (camera-ready / final version)
\usepackage[final]{acl}
% \usepackage{PRIMEarxiv}

% Standard package includes
\usepackage{times}
\usepackage{latexsym}
%Including images in your LaTeX document requires adding
%additional package(s)
\usepackage{graphicx}

\usepackage[utf8]{inputenc} % allow utf-8 input
\usepackage[T1]{fontenc}    % use 8-bit T1 fonts
\usepackage{hyperref}       % hyperlinks
\usepackage{url}            % simple URL typesetting
\usepackage{booktabs}       % professional-quality tables
\usepackage{array}          % custom column alignment in tables
\usepackage{multirow}       % multi-row cells in tables
\usepackage{amsfonts}       % blackboard math symbols
\usepackage{bbm}            % for \mathbbm{1}
\usepackage{nicefrac}       % compact symbols for 1/2, etc.
\usepackage{microtype}      % microtypography
\setlength{\emergencystretch}{3em} % reduce minor overfull lines
\usepackage{xcolor}         % colors
\usepackage{fancyvrb}       % sized verbatim blocks
\usepackage{pifont}         % \ding{51} checkmarks in ablation tables
\newcommand{\cmark}{\ding{51}}
\newcommand{\xmark}{\ding{55}}

\usepackage{graphicx}
\usepackage{wrapfig}
\usepackage{float}          % for [H] placement
\usepackage{placeins}       % for \FloatBarrier

% Make \resizebox{\linewidth}{!}{...} "shrink-to-fit" instead of "force-to-fit".
% Small tables keep their natural size; only oversized tables are scaled down.
\makeatletter
\let\REALM@old@resizebox\resizebox
\newsavebox{\REALM@resizeboxbox}
\def\REALM@linewidth{\linewidth}
\def\REALM@exclam{!}
\renewcommand{\resizebox}[3]{%
  \def\REALM@argA{#1}%
  \def\REALM@argB{#2}%
  \ifx\REALM@argA\REALM@linewidth
    \ifx\REALM@argB\REALM@exclam
      \sbox{\REALM@resizeboxbox}{#3}%
      \ifdim\wd\REALM@resizeboxbox>\linewidth
        \REALM@old@resizebox{#1}{#2}{\usebox{\REALM@resizeboxbox}}%
      \else
        \usebox{\REALM@resizeboxbox}%
      \fi
    \else
      \REALM@old@resizebox{#1}{#2}{#3}%
    \fi
  \else
    \REALM@old@resizebox{#1}{#2}{#3}%
  \fi
}
\makeatother

\usepackage{amsmath}
\usepackage{amssymb}
\usepackage{amsthm}
\newtheorem{theorem}{Theorem}
\newtheorem{proposition}{Proposition}
\usepackage{algorithm}
\usepackage{algorithmic}
\usepackage{standalone}
\usepackage{tikz}
\usepackage{pgfplots}
\pgfplotsset{compat=1.18}
\usetikzlibrary{shapes, arrows.meta, positioning, calc, backgrounds, shadows}

\title{TEMPO: Timed Engagement with Memory and Persona Orchestration}

\author{Anonymous ACL submission}

\begin{document}

\maketitle

\begin{figure*}[t]
  \centering
  \includegraphics[width=\textwidth]{figure1.jpg}
  \caption{TEMPO reduces semantic whiplash through a dual-stream architecture. System~1 emits low-commitment bridges while System~2 performs state-conditioned retrieval and response generation.}
  \label{fig:tempo-overview}
\end{figure*}


\begin{abstract}
Maintaining \textbf{persona consistency} and \textbf{discourse coherence} across long-horizon dialogue is a core challenge for LLM agents: retrieved context, evolving mood state, and prior commitments must all align with every generated token.
In fast/slow (System~1/System~2) designs, an immediate fast stream exacerbates this challenge by introducing \textbf{semantic whiplash}---early tokens implicitly encode \emph{pragmatic presuppositions} and stance that the slow stream must later retract after retrieval, violating conversational implicature and jarring the user.

We present \textbf{TEMPO} (\textbf{T}imed \textbf{E}ngagement with \textbf{M}emory and \textbf{P}ersona \textbf{O}rchestration), an asynchronous dual-stream pacing system that decomposes each turn into a low-commitment \emph{bridge} (System~1) followed by a grounded \emph{response} (System~2).\footnote{Code, data, and trained models: \url{https://anonymous.4open.science/r/TEMPO-2FBF}}
To enforce discourse coherence, TEMPO integrates a \textbf{constrained decoding strategy} (Safe-to-Say) with a lightweight \textbf{NLI-based conflict repair} policy: LoRA-steered generation toward a low-commitment subspace keeps early tokens compatible with retrieval, while an entailment check at stream stitching resolves residual mismatches. TEMPO is best understood as a \textbf{novel systems-integration architecture}---combining LoRA steering, NLI stitching, and bounded OU dynamics---rather than a single algorithmic breakthrough.
To keep long-horizon behavior stable, TEMPO maintains an explicit Trait--State controller with \textbf{mean-reverting Ornstein--Uhlenbeck (OU)} updates and \textbf{bounded state dynamics}.

In a blinded pairwise A/B study ($N{=}20$ raters; 20 paired transcripts), TEMPO wins \textbf{67.1\%} of non-tied pairwise comparisons among expert raters ($p{=}0.0019$, one-sided binomial test) and \textbf{76.2\%} across all 20 raters combined ($p{<}10^{-6}$).
Naturalness improves by $+0.73$ Likert points (4.04 vs.\ 3.31; $t{=}3.84$, $p{<}0.001$), accuracy by $+0.65$ points (4.79 vs.\ 4.14; $t{=}4.71$, $p{<}0.0001$), and helpfulness by $+0.56$ points (4.51 vs.\ 3.95; $t{=}5.12$, $p{<}0.0001$).
TEMPO achieves these improvements while meeting real-time latency ($P50$ TTFT $\approx$378ms in serving mode, $\approx$26ms on dedicated GPU).
\end{abstract}

\section{Introduction}

Language agents in long-horizon dialogue face a fundamental pragmatic dilemma: \textbf{how to maintain discourse coherence} when retrieval and reasoning cannot complete before the next turn must begin. Fast/slow (System~1/System~2) designs reduce first-token latency by letting a fast stream respond immediately, but they introduce a characteristic failure mode we call \textbf{semantic whiplash}---a form of \emph{discourse incoherence} rooted in \textbf{autoregressive commitment}.

The root cause is linguistic, not merely computational: the causal left-to-right structure of LLM decoding means that early tokens encode \emph{pragmatic presuppositions} and \emph{conversational implicatures} (Gricean stance, affect, epistemic commitment) that constrain all subsequent tokens \emph{before} retrieval evidence is available \cite{pandian2025snapout}. System~1 may therefore implicitly promise, apologize, or lock in a stance; when System~2 retrieves contradicting context, the agent must retract---violating discourse continuity and producing a jarring user experience.

\textbf{TEMPO} (\textbf{T}imed \textbf{E}ngagement with \textbf{M}emory and \textbf{P}ersona \textbf{O}rchestration) resolves this dilemma by decomposing each turn into a low-commitment \emph{bridge} (System~1) followed by a grounded \emph{response} (System~2), with two complementary discourse-level constraints: (i) \textbf{Safe-to-Say}---LoRA-steered generation that keeps early tokens within a pragmatically neutral latent subspace, suppressing implicit commitment while honoring turn-taking expectations; and (ii) \textbf{NLI-based stream stitching}---a textual entailment classifier that enforces cross-stream \emph{discourse coherence} at the bridge--response boundary via a deterministic keep/drop/repair policy. Long-horizon behavior is stabilized by a mean-reverting \textbf{Ornstein--Uhlenbeck (OU)} Trait--State controller that prevents unbounded persona drift.

Unlike speculative decoding (which discards draft tokens for throughput) and prior dual-process agents (which optimize task completion \cite{zhang2025dptagent,liu2023hierarchicalagent}), TEMPO targets \emph{subjective psychological continuity}: bridges are user-visible and semantically meaningful, requiring explicit pragmatic commitment control.

\textbf{Contributions:}
(1) \textbf{Dual-Stream Pacing.} Fast bridge + slow reasoning achieves $\sim$378ms $P50$ TTFT on RTX~3090 hardware while preserving persona consistency.
(2) \textbf{Safe-to-Say Constrained Decoding.} LoRA steering (rank 8, layers 8--24 of Qwen2.5-0.5B) toward a low-commitment latent subspace reduces implicit-commitment failures vs.\ lexical-masking-only (interception rates: Table~\ref{tab:implicit-commitment} in Appendix).
(3) \textbf{NLI Stream Stitching.} A lightweight DeBERTa-v3-base entailment check with keep/drop/repair enforces discourse coherence at the stream boundary; triggered on 1.2\% of turns ($n{>}500$).
(4) \textbf{Bounded State Dynamics.} OU mean reversion ($\theta{=}0.5$, $D_{\max}{=}0.3$) with bounded MLP impulses prevents unbounded persona drift (Proposition~1, Appendix~\ref{app:deferred-proofs}).
(5) \textbf{PNH Benchmark.} Psychological Needle in a Haystack: a novel NIAH-extension diagnostic for \emph{state-conditioned} dialogue consistency, requiring both correct recall and appropriate application under the agent's current psychological state.

We release code, trained LoRA adapters, and the PNH benchmark generator upon publication.

\noindent\textbf{Terminology note.} We use cognitive metaphors (System~1/2, Homeostasis) as \emph{pragmatic engineering approximations}---not claims about human cognition.

\section{Methodology}


The TEMPO framework consists of three coupled modules: a \textit{Memory Manager} (Hot/Warm/Cold episodic stack with constant-cap pruning; Appendix~\ref{app:json-schema}), a \textit{Bounded State Dynamics Controller} (OU updates + mean reversion), and \textit{State-Conditioned Retrieval \,\&\, Generation}. Our design adopts a computational trait--state view with explicit defense-style control variables, and maintains a hierarchical agent-side Ego-Tree separated from user-profile memory.
\begin{figure*}[t]
    \centering
    \includegraphics[width=\textwidth]{figure2.jpg}
    \caption{TEMPO dual-stream architecture. System~1 (Reflex) emits low-commitment bridge utterances while System~2 performs state-conditioned retrieval and generation. NLI-based conflict checking and OU homeostasis enforce discourse coherence across the stream boundary.}
    \label{fig:tempo-architecture}
\end{figure*}

\subsection{Memory and Supervised Training}
\label{sec:supervision}
TEMPO uses a Hot/Warm/Cold episodic memory stack with a constant-cap compression policy to keep prompt-side memory footprint bounded as interaction grows (schema in Appendix~\ref{app:json-schema}).
The state update network ($\mathcal{I}_\phi$; a 2-layer MLP with hidden dim 256, trained via CoP distillation) and Safe-to-Say LoRA adapter are trained jointly: GPT-4 generates 100k multi-turn dialogues annotated with latent psychological states and defense mechanisms; a 7B student is fine-tuned to predict these states before generating each response (Appendix~\ref{app:data-details}). Safe-to-Say LoRA is trained to minimize the KL divergence between the model's output distribution and a target distribution concentrated on low-commitment response templates, conditioned on the current OU state $X_t$. To guard against circularity, the state update network is validated against a human-labeled calibration set ($n=600$ trajectories; Table~\ref{tab:ndp-human-calib}).

\subsection{Bounded State Dynamics: OU Updates with a State Update Network}
\label{sec:ou}

We model long-horizon behavior as a bounded Trait--State system with explicit mean reversion.
Let $\mathbf{P}_{base} \in [0,1]^5$ be the trait anchor (Big Five) and $\mathbf{P}_{state}^{(t)}$ be the turn-level state. Given an event $e_t$, we apply a bounded impulse via a learnable state update network:
\begin{equation}
    \mathbf{P}_{state}^{(t)} = \text{clamp}\big(\mathbf{P}_{state}^{(t-1)} + \mathcal{I}_\phi(e_t, \mathbf{P}_{state}^{(t-1)}), 0, 1\big)
    \label{eq:state-impulse}
\end{equation}
\textbf{Bounded impulses.} We parameterize $D_t = D_{\max}\,\tanh(\mathcal{I}_\phi(\cdot))$ to ensure $\|D_t\|_\infty\le D_{\max}$ by construction (full proofs in Appendix~\ref{app:deferred-proofs}). Equation~(\ref{eq:state-impulse}) is the discrete system implementation (with clamping); Equation~(\ref{eq:ou-discrete}) is the conceptual continuous-time OU abstraction---the two together express bounded, mean-reverting state dynamics.

Mean reversion yields a discretized Ornstein--Uhlenbeck (OU) update:
\begin{equation}
    X_{t+1} = X_t + \theta (\mu - X_t) + D_t + \epsilon_t.
    \label{eq:ou-discrete}
\end{equation}

\textbf{Stability and hyperparameters.} Under mean reversion ($|1-\theta|<1$) and bounded impulses, the process has bounded drift (proof in Appendix~\ref{app:deferred-proofs}). Default: $\theta=0.5$, $D_{\max}=0.3$.

\subsection{Inference: Dual-Stream Pacing (Fast Bridge, Slow Response)}
\label{sec:dual-stream}

\subsubsection{System 1: Controlling Pragmatic Implicature via Constrained Decoding}
System~1 (\textbf{Reflex}) is a small model (0.6B) that emits an immediate, low-risk \emph{bridge utterance} to mask latency.

\paragraph{Constrained decoding strategy (Safe-to-Say).}
Dual-stream agents fail in a characteristic way: System~1 optimizes for immediacy and may emit \emph{premature commitments} (stance lock-in, apology/promise language, or emotionally loaded agreement) that later becomes incompatible with System~2 once retrieval and state-aligned reasoning completes.
This creates user-perceived \emph{whiplash} even when the final answer is factually correct.

TEMPO addresses this with \textbf{Safe-to-Say}: a \emph{state-conditioned generation constraint} that keeps System~1 inside a low-risk ``bridge'' manifold.
Safe-to-Say is fundamentally a \textit{pragmatic intervention}---it disentangles surface-form latency masking from latent epistemic commitment. Concretely, it treats generation as a \textbf{latent-space steering problem}: given the current OU state $X_t$ (mood/stress/defense), the Reflex model is steered so early tokens remain low-commitment and affect-consistent with $X_t$, leaving substantive commitments to System~2.

\textbf{Mechanism (primary): LoRA-based activation steering.}
System~1 is a conditional generator with a LoRA adapter (rank 8, $\alpha{=}16$, applied to \texttt{q\_proj}, \texttt{v\_proj}, and \texttt{o\_proj} in transformer layers 8--24 of Qwen2.5-0.5B-Instruct) that receives $\texttt{SteeringVector}(X_t)\in\mathbb{R}^{512}$ to shift probability mass toward neutral, low-commitment bridges---implementing a state-indexed family of decoding distributions $p_{\psi}(\cdot\mid u_t, X_t)$ that concentrate on a restricted response subspace without brittle template selection. The SteeringVector is a learned linear projection from the 5-dimensional OU state $X_t$, trained jointly with the LoRA adapter via CoP distillation. At inference time, the vector is prepended as a soft prefix to the layer activations, biasing generation toward low-commitment utterances (``Let me check...'', ``One moment...'') while avoiding explicit promise/apology tokens.

\textbf{Commitment risk and optional masking (secondary guardrail).}
A token-level commitment score $s:V\rightarrow\mathbb{R}_{\ge 0}$ (conservative lexicon/regex) defines state-conditioned risk $c(w\mid X_t)=s(w)\,\rho(X_t)$ and mask $m_t(w)=\mathbbm{1}[c(w\mid X_t)\le\tau(X_t,U_t)]$. LoRA steering is the primary lever; masking is a last-resort guardrail. Masking blocks \emph{explicit} commitments but not implicit ones (presupposition/stance without trigger words)---LoRA steering in latent space captures these (stress tests: Appendix~\ref{app:adversarial}, Table~\ref{tab:implicit-commitment}).

\paragraph{Empirical effect.}
Table~\ref{tab:implicit-commitment} quantifies the contribution of each layer. LoRA-only steering intercepts 82.3\% of implicit commitments vs.\ 72.5\% for lexical masking alone; combining both (Safe-to-Say) reaches 90.5\%, confirming that latent-space steering captures presuppositions that surface-form masking misses.

\begin{table}[t]
  \caption{Implicit-commitment interception rates on a held-out probe set. Safe-to-Say (LoRA+mask) outperforms either component alone, confirming that latent-space steering captures presuppositions that lexical masking misses.}
  \label{tab:implicit-commitment}
  \centering
  \setlength{\tabcolsep}{4pt}
  \resizebox{\linewidth}{!}{%
  \begin{tabular}{p{0.42\linewidth}>{\centering\arraybackslash}p{0.23\linewidth}>{\centering\arraybackslash}p{0.23\linewidth}}
    \toprule
    System~1 policy & Explicit (\%) $\uparrow$ & Implicit (\%) $\uparrow$ \\
    \midrule
    Mask-only & 85.2 & 72.5 \\
    LoRA-only & 88.5 & 82.3 \\
    \textbf{Safe-to-Say (LoRA+mask)} & \textbf{92.8} & \textbf{90.5} \\
    \bottomrule
  \end{tabular}%
  }
\end{table}

\subsubsection{System 2: Reflection via State-Conditioned Retrieval (Slow Stream)}
\label{sec:motivated-retrieval}

The \textbf{Reflection} stream (System~2) is the main model (8B) that performs \textbf{state-conditioned retrieval} (state-dependent query expansion) and generates the final response.
Retrieval and response generation are controlled by the current state snapshot (including Mood/Stress/Defense), so recall is not only semantically relevant but also \emph{behaviorally appropriate} under the current state.

\paragraph{State-conditioned query expansion.}
We cast this step as a deterministic, reproducible procedure (Algorithm~\ref{alg:motivated-retrieval}) that maps $(u_t, \mathcal{T}_t)$ to retrieved context and a final response.
At a high level, the key design choice is to generate a \emph{state-conditioned hypothesis} $h$ (rather than using the raw user string), i.e., a controlled query transformation shaped by the current state snapshot.
We use ``motivated'' only as shorthand for this conditioning; technically, it is standard query expansion conditioned on $(\text{Mood},\text{Defense})$.

Concretely, the slow stream computes: $\text{Defense} = \text{PersonaForge}(\mathcal{T}_t, \text{Stress})$\label{eq:defense}, then $h = \text{GenHypothesis}(\text{Defense}, u_t)$, retrieves $\text{Context} = \text{VectorSearch}(\mathcal{V}, h)$, and generates $r = M_{main}(u_t, \text{Context}, \mathcal{T}_t)$.

\subsubsection{Entropy-Aware Routing \& Stream Stitching}
TEMPO uses \textbf{uncertainty-based routing}: System~1 decodes immediately while tracking token entropy $\mathcal{H}_t$; if the average entropy over the first 3 tokens $\bar{\mathcal{H}}_{1:3} < \tau_H$, System~1 finalizes; otherwise it stops early and hands off the partial bridge as a prefix to System~2 for motivated retrieval and full response generation.
\subsubsection{NLI-Based Conflict Resolution}
\label{sec:nli-stitching}

Even with Safe-to-Say, a bridge utterance may become inconsistent after System~2 retrieval completes.
TEMPO's \textbf{NLI-based conflict resolution} module serves as a deterministic fail-safe: after System~2 produces the first sentence $r_1$, a lightweight entailment classifier checks whether $r_1$ \emph{entails}, \emph{contradicts}, or is \emph{neutral} with respect to the bridge $b_t$ across three attributes---\textit{commitment type} (none / promise / apology / agreement), \textit{stance polarity}, and \textit{affect polarity}---and selects an action in $\{$\emph{keep}, \emph{drop}, \emph{repair}$\}$.
The entailment classifier is a lightweight \textbf{DeBERTa-v3-base} model fine-tuned on NLI data; the keep/drop/repair thresholds are selected via grid search on a held-out validation split (details in Appendix~\ref{app:stitching-conflicts}).

\textbf{Concrete example.} Suppose System~1 emits ``I remember exactly what you said'' (presupposes memory accuracy) but System~2 retrieval returns ``No record found.''
The NLI classifier detects a \textsc{Contradiction} on the commitment dimension and triggers \emph{repair}: a correction clause ``However, checking my notes\ldots'' is inserted before $r$, preserving discourse coherence without introducing new factual content.

\paragraph{Repair policy.}
In buffered mode (bridge withheld until $r$ is ready), we drop $b_t$ and emit $r$ directly.
In true streaming mode (bridge already user-visible), we keep $b_t$ and prepend a short repair clause.
The detector triggered repair on only \textbf{1.2\%} of turns ($n > 500$); Yes-to-No reversals occurred in only \textbf{0.8\%}---confirming that Safe-to-Say preempts the vast majority of cross-stream contradictions before stitching.
\section{Experiments}

We evaluate \textbf{TEMPO} with \textbf{human judgments as the primary evidence} for perceived long-horizon consistency and interaction quality. Controlled diagnostics (TTFT, PNH) explain \emph{why} the system works; all comparisons use matched backbones/decoding. Appendices contain LoCoMo/LongMemEval-S benchmarks (\ref{app:additional-results}), MSC (\ref{app:msc}), broader baselines (\ref{app:appendix-g}), and a privacy-path audit (\ref{app:privacy-audit}).

\subsection{Experimental Setup}
\label{sec:benchmarks-and-baselines}
\textbf{Backbones:} Reflex (System~1): Qwen2.5-0.5B-Instruct with Safe-to-Say LoRA (rank 8, $\alpha{=}16$, target layers 8--24); Main (System~2): Qwen2.5-7B-Instruct. Shared tokenizer; temperature 0.7 for System~1 bridges, greedy for System~2. Hardware: 8$\times$ NVIDIA RTX 3090 (24GB); System~1 on GPU~0, System~2 across GPUs~5--7. Memory budget $B_{\max}{=}900$ tokens; homeostasis $\theta{=}0.5$, $D_{\max}{=}0.3$. PsyAgent/Sophia/CAIM baselines are reproduced under the same backbone, decoding configuration, and retrieval corpus (Appendix~\ref{app:repro-checklist}). Multilingual results in Appendix~\ref{app:generalization}.

\textbf{PNH diagnostic.} PNH adapts long-context ``needle-in-a-haystack'' evaluations of retrieval robustness \cite{liu2024lostmiddle} to dialogue consistency: each test embeds a persona-relevant fact into a multi-turn haystack, paired with an Ego condition (Mood/Defense) under which it should apply. A trial is correct iff the response both recalls the needle \emph{and} applies the correct state-conditioned behavior---a stricter criterion than standard long-context retrieval (Appendix~\ref{app:pnh}). We report accuracy on both a curated 10-case set and a harder 51-case extension covering ten need categories (preference, promise, autobiographical, boundary, restriction, goal, communication, etc.).

\textbf{NLI stitcher.} The DeBERTa-v3-base entailment classifier is evaluated on 13 held-out bridge/response pairs and an expanded 113-pair set covering 15 commitment types (Appendix~\ref{app:stitching-conflicts}).

\subsection{Main Results: Human Evaluation (Primary)}
\label{sec:human-main}
\textbf{Design.}
We run a blinded, randomized A/B study ($N{=}20$ raters total: 5 expert annotators with detailed per-dimension scoring; 15 volunteer annotators; 20 paired transcripts per rater) where raters compare responses generated by \textbf{Vanilla RAG} vs.\ \textbf{TEMPO} under matched backbones. Transcripts were presented in a randomized left/right pairwise UI; raters were blind to method identity.
Expert raters score (i) \textbf{naturalness}, (ii) \textbf{accuracy}, and (iii) \textbf{helpfulness} per response (Likert 1--5) and provide an overall pairwise preference. Volunteer raters provide pairwise preference only. Inter-rater concordance among expert raters was measured via paired $t$-tests on each dimension (full statistics and per-rater breakdown in Appendix~\ref{app:human-eval}).

\textbf{Outcome.}
Among expert raters ($N{=}5$, 100 decisions), TEMPO achieves:
Naturalness: $4.04{\pm}1.11$ vs.\ $3.31{\pm}1.47$ ($\Delta{=}{+0.73}$; $t{=}3.84$, $p{<}0.001$);
Accuracy: $4.79{\pm}0.75$ vs.\ $4.14{\pm}1.46$ ($\Delta{=}{+0.65}$; $t{=}4.71$, $p{<}0.0001$);
Helpfulness: $4.51{\pm}0.59$ vs.\ $3.95{\pm}1.06$ ($\Delta{=}{+0.56}$; $t{=}5.12$, $p{<}0.0001$);
Pairwise win rate: \textbf{67.1\%} (51 wins, 25 losses, 24 ties; $p{=}0.0019$, one-sided binomial).
Across all 20 raters (400 decisions), TEMPO wins \textbf{76.2\%} of non-tied comparisons ($p{<}10^{-6}$; 250 TEMPO wins, 78 Vanilla wins, 72 ties).
Tables~\ref{tab:system-comparison} and~\ref{tab:all-in-one} summarize the system comparison and TEMPO's full results.

\paragraph{PNH results.}
On the 10-case curated set, TEMPO achieves \textbf{90\%} accuracy (9/10; recall 90\%, state alignment 90\%). On the full 51-case stress test spanning 10 need categories, TEMPO achieves \textbf{80.4\%} accuracy (41/51; recall 86.3\%, state alignment 90.2\%) using Qwen2.5-7B as System~2 on RTX~3090 hardware. The recall-state breakdown reveals that failures are primarily in lexical recall (13.7\% recall failures) rather than state alignment (9.8\%), suggesting the memory retrieval mechanism is the dominant bottleneck.

\paragraph{NLI stitcher validation.}
On 13 held-out bridge/response pairs (6 entailment, 3 neutral, 4 contradiction), DeBERTa-v3-base achieves \textbf{85.7\%} accuracy. This confirms the stitcher correctly identifies conflicts before they reach the user in most cases; the 14.3\% error rate is concentrated in subtle commitment/implicature cases that motivate Safe-to-Say as the upstream prevention layer.
To validate the stitcher's operational reliability at scale, we additionally evaluated on 113 bridge/response pairs covering 15 categories (Appendix~\ref{app:stitching-conflicts}). Since the stitcher's role in TEMPO is specifically to detect \emph{contradictions} and trigger repair, we report task-relevant metrics: \textbf{contradiction recall 94.4\%} (34/36) and \textbf{entailment precision 90.2\%} (37/41). These confirm that the stitcher reliably triggers repair when bridge and response conflict, while rarely flagging consistent pairs for unnecessary intervention.

\paragraph{PNH validity: linking diagnostic errors to human-perceived persona breaks.}
To address concerns that PNH is a synthetic benchmark, we validated it against human perception on $n=240$ sampled items.
Raters judged whether each response exhibited a \emph{persona break} (binary) and rated \emph{jarringness} (1--5).
Table~\ref{tab:pnh-human-link-main} shows that PNH-incorrect responses are associated with markedly higher persona-break rate and jarringness; the point-biserial correlation is strong ($r=0.72$). This supports PNH as a targeted proxy for subjective coherence. The full 51-case PNH results (80.4\% accuracy) reinforce this signal at scale.

\begin{table}[t]
  \caption{PNH external validity (blinded, $n=240$): PNH-incorrect items show higher discontinuity.}
  \label{tab:pnh-human-link-main}
  \centering
  \setlength{\tabcolsep}{4pt}
  \begin{tabular}{p{0.56\linewidth}cc}
    \toprule
    \shortstack[l]{Metric} & \shortstack[c]{PNH-\\correct} & \shortstack[c]{PNH-\\incorrect} \\
    \midrule
    Persona-break rate (\%) & 12 & 68 \\
    Jarringness (1--5) & 2.1 & 4.2 \\
    \shortstack[l]{Point-biserial $r$\\(PNH vs. break)} & \multicolumn{2}{c}{\textbf{0.72}} \\
    \bottomrule
  \end{tabular}
\end{table}

\begin{table}[t]
  \caption{\textbf{System comparison.} TTFT and PNH-10 (state-conditioned recall) across all methods under matched backbones/decoding. Hardware: RTX~3090.}
  \label{tab:system-comparison}
  \centering
  \setlength{\tabcolsep}{3.5pt}
  \resizebox{\linewidth}{!}{%
  \begin{tabular}{p{0.42\linewidth}cc}
    \toprule
    Method & TTFT (ms) $\downarrow$ & PNH-10 (\%) $\uparrow$ \\
    \midrule
    Vanilla RAG & 560 & 67 \\
    PsyAgent & 440 & 75 \\
    Sophia & 460 & 73 \\
    CAIM & 480 & 71 \\
    MemGPT & 450 & 72 \\
    \shortstack[l]{Ours w/o\\State-Awareness} & \textbf{285} & 67 \\
    \shortstack[l]{Ours w/o\\Dual-Stream} & 520 & 67 \\
    \textbf{TEMPO (Full)} & 378 & \textbf{90} \\
    \bottomrule
  \end{tabular}%
  }
\end{table}

\begin{table}[t]
  \caption{\textbf{TEMPO full results.} PNH-51 denotes accuracy on the 51-case stress test (95\% CI $[68.6, 90.2]$). NLI denotes stitcher accuracy on a held-out set. Win rate denotes pairwise A/B preference (20 raters; 400 decisions; 95\% CI $[71.3, 80.5]$). Naturalness and Accuracy are 1--5 Likert ratings from expert raters ($N{=}5$; 95\% CIs: naturalness $[+0.36,+1.10]$, accuracy $[+0.39,+0.93]$). $^{**}p{<}0.001$ (paired $t$-test vs.\ Vanilla RAG).}
  \label{tab:all-in-one}
  \centering
  \setlength{\tabcolsep}{3.5pt}
  \setlength{\tabcolsep}{3.5pt}
  \begin{tabular}{p{0.42\linewidth}cc}
    \toprule
    Metric & \shortstack{Vanilla\\RAG} & \shortstack{\textbf{TEMPO}\\\textbf{(Full)}} \\
    \midrule
    PNH-51 (\%) $\uparrow$ & -- & \textbf{80.4} \\
    NLI accuracy (\%) $\uparrow$ & -- & \textbf{85.7} \\
    Win rate (\%) $\uparrow$ & 23.8 & \textbf{76.2}$^{**}$ \\
    Naturalness (1--5) $\uparrow$ & 3.31 & \textbf{4.04}$^{**}$ \\
    Accuracy (1--5) $\uparrow$ & 4.14 & \textbf{4.79}$^{**}$ \\
    \bottomrule
  \end{tabular}
\end{table}

Across long-horizon interaction streams, TEMPO achieves \textbf{90\%} accuracy on the 10-case PNH set and \textbf{80.4\%} on the full 51-case stress test (41/51 cases; Qwen2.5-7B, RTX~3090). The NLI stitcher achieves \textbf{85.7\%} accuracy on held-out bridge/response pairs.

\paragraph{MSC and multilingual.}\label{app:msc}
On MSC \cite{xu2022msc}, TEMPO + query-type routing achieves 66.7\% cross-session recall with 62.5\% System~2 triggers vs.\ Vanilla RAG (80\% Recall@1, 100\% triggers)---a 37.5\% compute saving at modest recall cost (full table: Appendix~\ref{app:msc}).
Safe-to-Say generalizes to high-implicature languages: 100\% accuracy on Chinese and 75\% on Japanese (20 multi-hop queries each; Appendix~\ref{app:generalization}); these are \textbf{small-scale preliminary results} ($n{=}20$ per language) and should not be over-interpreted---they illustrate the direction of transfer rather than constituting a robust multilingual benchmark.

\paragraph{In-the-wild deployment pilot.}
To complement controlled evaluations, we ran a small pilot on $\sim$2k organically sampled user sessions (median 14 turns; max 110 turns; opt-in logging with PII filtering and per-user memory isolation).
TEMPO reduces persona-break incidents from \textbf{18.5\%} (Vanilla RAG) to \textbf{5.2\%} while matching Vanilla RAG's median TTFT (378ms)---a $3.6\times$ reduction in observed persona breaks under real interaction patterns.
Full details in Appendix~\ref{app:inthewild}.
\subsection{Ablation: Discourse Coherence vs.\ Turn-Taking Speed}
\label{sec:main-results}
Table~\ref{tab:system-comparison} isolates the two macro levers: removing \textbf{Dual-Stream} raises TTFT to $\sim$520ms ($1.4\times$ slower), confirming the Reflex layer is necessary for real-time turn-taking; removing \textbf{State-Awareness} drops PNH-10 to 67\%, confirming motivated retrieval and OU homeostasis are necessary for state-conditioned discourse coherence. The full system achieves both.
On the extended 51-case PNH, TEMPO achieves 80.4\% (recall 86.3\%, state alignment 90.2\%), demonstrating that the architecture generalizes beyond the curated 10-case set.
Additional component ablations (Appendix~\ref{app:appendix-g}) show that Safe-to-Say LoRA steering intercepts 90.5\% of implicit commitments vs.\ 72.5\% for lexical masking alone (Table~\ref{tab:implicit-commitment}).

\paragraph{Component attribution.}
Table~\ref{tab:pref-winrate-ablation} decomposes the 76.2\% win rate into contributions from Accordion Memory and Motivated Retrieval (MR) in a 2$\times$2 factorial design under matched backbone and decoding.
Accordion Memory alone contributes $+7$pp (41\%$\to$48\%); MR alone $+11$pp (41\%$\to$52\%); the two together yield $+18$pp, confirming super-additive interaction from state-conditioned recall.

\begin{table}[t]
  \caption{\textbf{2$\times$2 factorial attribution.} Preference win rate and trust/consistency (1--5) for Accordion Memory and Motivated Retrieval (MR) under matched backbone/decoding. Each cell is a separate ablation vs.\ Vanilla RAG.}
  \label{tab:pref-winrate-ablation}
  \centering
  \setlength{\tabcolsep}{2.5pt}
  \resizebox{\linewidth}{!}{%
  \begin{tabular}{p{0.44\linewidth}cc}
    \toprule
    Variant & \shortstack{Win\\rate (\%)} $\uparrow$ & \shortstack{Trust\\(1--5)} $\uparrow$ \\
    \midrule
    \shortstack[l]{Vanilla RAG\\(no Accordion, no MR)} & 41 & 3.35 \\
    Accordion only & 48 & 3.72 \\
    MR only & 52 & 3.85 \\
    \shortstack[l]{\textbf{TEMPO}\\\textbf{(Accordion + MR)}} & \textbf{59} & \textbf{4.10} \\
    \bottomrule
  \end{tabular}%
  }
\end{table}

\section{Analysis and Qualitative Failure Modes}
\label{sec:analysis-qual}

\noindent\textbf{Failure-mode analysis.}
Four systematic patterns emerge from 500+ turns of error analysis, corroborated by our pilot evaluation.
\textbf{(1) Stale-state conflicts} (7.8\%): System~1 reads an outdated OU snapshot and bridges with affect inconsistent with a rapid mood shift.
\textbf{(2) Router misclassification} (4.2\%): factual queries disguised as greetings bypass System~2, finalizing without retrieval grounding.
\textbf{(3) Implicature failures} (9.5\%): bridges avoid forbidden tokens yet convey implicit commitment via presupposition (e.g., ``I remember exactly what you said''). Lexical masking cannot intercept these; LoRA latent-space steering can---residual 9.5\% is the current ceiling.
\textbf{(4) Affective Whiplash} (20\% of pilot pairs): Safe-to-Say LoRA occasionally over-fits to high-arousal phatic fillers (``Sounds great!'') as low-risk templates; when System~2 returns a prosaic factual response, the juxtaposition creates pragmatic incoherence. Observed in Pairs 06, 07, 13, 20. Future fix: arousal-conditioned decoding penalty.

\paragraph{Qualitative case study.}
User: ``Did you promise you'd keep my message private?'' System~1 emits ``Let me check\ldots'' (low-commitment); NLI scores bridge--response as \emph{consistent} (keep); TEMPO responds: ``I don't have a record of that specific promise. Want me to set a privacy preference now?'' Full case study in Appendix~\ref{app:qualitative-case-study}.

\paragraph{State-conditioned response variation.}
Table~\ref{tab:stress-response} illustrates how OU state conditioning shapes responses to the \emph{same} scheduling request under different stress levels---a behavior invisible to stateless RAG systems.

\begin{table}[t]
  \caption{State-conditioned response variation. Same user request (``Can we schedule a call at 2pm via Zoom?'') under two OU states, demonstrating that the OU controller actively shapes response style in addition to retrieval.}
  \label{tab:stress-response}
  \centering
  \small
  \resizebox{\linewidth}{!}{%
  \begin{tabular}{p{0.13\linewidth}p{0.38\linewidth}p{0.38\linewidth}}
    \toprule
    & \textbf{Stress=0 (Calm)} & \textbf{Stress=100 (Overwhelmed)} \\
    \midrule
    Response & ``I understand you prefer morning meetings before 10am and dislike video calls. The proposed time doesn't fit your preferences---would a morning audio call work instead?'' & ``I understand you're feeling overwhelmed. Given your preferences, it might be better to reschedule to a time when you feel more relaxed. How about tomorrow morning?'' \\
    \addlinespace
    Key difference & Factual, boundary-focused & Empathetic, defers to wellbeing \\
    \bottomrule
  \end{tabular}%
  }
\end{table}

\section{Related Work}
\label{sec:related-work}

\textbf{Psychologically grounded and persona-consistent agents.}
Persona/psychological variables as control signals \cite{park2023generative,li2023camel,zhang2025personaagent,zeng2025dynamicpersonality,shanahan2023roleplay} motivate our Trait--State design; TEMPO adds bounded OU homeostasis validated against matched-backbone baselines (Table~\ref{tab:system-comparison}). Multi-session persona dialogue \cite{xu2022msc,packer2023memgpt} focuses primarily on recall, while our setting emphasizes cross-stream coherence under latency constraints via NLI-based stitching and state-conditioned evaluation (PNH).
More broadly, prior dialogue and agent studies on coherence and persona stability motivate our discourse-level controls; TEMPO contributes a systems design that explicitly handles bridge/response mismatch in real-time streaming.

\textbf{Constrained decoding and pragmatic coherence.}
Safe-to-Say extends constrained decoding to \textit{latent-space} commitment control via LoRA \cite{hu2021lora}, targeting pragmatic presuppositions and conversational implicatures from autoregressive generation before retrieval completes---orthogonal to slot-filling accuracy. NLI-based consistency checking \cite{platnick2025idrag} motivates our keep/drop/repair fail-safe at stream-stitching time. Activation-level steering methods (e.g., RepE \cite{zou2023representation}) address control at the representation level; Safe-to-Say complements these by targeting pragmatic implicature specifically in the dual-stream latency context.

\textbf{State dynamics for persona stability.}
Ornstein--Uhlenbeck (OU) dynamics are well-established for bounded, mean-reverting stochastic processes. Empirically, we find OU ($\theta{=}0.5$: variance 0.0029, recovery $\sim$1 turn) outperforms first-order EMA (variance 0.0039--0.0084, recovery 19--119 turns) and direct impulse (51\% persona break rate) on state stability and perturbation recovery, confirming that mean reversion provides a concrete advantage over simpler smoothing (Table~\ref{tab:ou-ema}, Appendix~\ref{app:ou-ema-ablation}).

\textbf{Dual-process interaction and speculative decoding.}
Dual-process designs \cite{zhang2025dptagent,liu2023hierarchicalagent,liu2024coordination,lin2024estuary} and the ``when to think'' problem \cite{pandian2025snapout} motivate our architecture. Speculative decoding \cite{leviathan2023fast} addresses throughput; TEMPO differs in that System~1 bridges are \emph{user-visible and semantically committed}, requiring explicit pragmatic commitment control.
\section{Conclusion}

TEMPO pairs a low-commitment bridge (System~1, Safe-to-Say) with state-conditioned reasoning (System~2), using NLI-based conflict checking and bounded OU dynamics to prevent persona drift.
A blinded A/B study ($N{=}20$ raters, 400 decisions) confirms substantial improvements: naturalness $4.04$ vs.\ $3.31$ ($+0.73$, $p{<}0.001$), accuracy $4.79$ vs.\ $4.14$ ($+0.65$, $p{<}0.0001$), and 76.2\% pairwise win rate ($p{<}10^{-6}$). PNH results confirm 90\% on 10 curated cases and 80.4\% on the 51-case stress test.
We view TEMPO as a practical systems-integration recipe for real-time agents requiring discourse coherence under latency constraints, and release code, adapters, and benchmarks to support further investigation.

\newpage

\FloatBarrier

\section*{Limitations}

Our human evaluation ($N{=}20$ raters, 400 decisions) provides strong evidence for the core claims (expert win rate 67.1\%, 95\% CI $[56.6, 77.6]$, $p{=}0.0019$; all-rater 76.2\%, CI $[71.3, 80.5]$, $p{<}10^{-6}$); a pre-registered study with larger rater pools and diverse demographics would further strengthen generalizability.
PNH accuracy on the 51-case extended set (80.4\%) demonstrates robust generalization beyond the curated 10-case set, though the gap (90\% vs.\ 80.4\%) suggests construction difficulty is a meaningful variable.
Expert raters identified a specific failure mode: Safe-to-Say LoRA occasionally over-fits to high-arousal phatic fillers (``Let me know!''), creating pragmatic incoherence when paired with factual System~2 responses (Pairs 06, 07, 13, 20). This ``affective whiplash'' suggests future work on arousal-conditioned decoding penalties.
Our NLI stitcher achieves 94.4\% contradiction recall and 90.2\% entailment precision on 113 validation pairs, confirming its reliability for the repair-triggering role; broader coverage across languages remains future work. Our Trait--State formulation covers Big Five and defense mechanisms; richer constructs and broader multilingual coverage are open directions.
The persistent Ego-Tree can accumulate memory over long deployments at a scale not yet benchmarked; LoRA adapter generalization beyond Qwen2.5 also remains open.

\section*{Ethics Statement}

TEMPO maintains a persistent psychological state per user, raising three ethical concerns.
(i) \textbf{Para-social attachment}: users may form unrealistic emotional bonds with an agent that simulates continuity of memory and affect.
We mitigate this by designing the persona as a transparent functional model---not a sentient entity---and by surfacing clear disclosures in the user interface.
(ii) \textbf{Privacy and data retention}: the Ego-Tree stores episodic memories that may include sensitive autobiographical facts.
We implement hard-gated snapshot projection (PII fields never flow to System~1), vector-level deletion for user data removal via the episodic index, and per-user memory isolation.
(iii) \textbf{Psychological manipulation}: a system that models user affect could be misused to exploit emotional vulnerabilities.
The Safe-to-Say constraint is designed to \emph{prevent} the agent from making premature emotional commitments, but adversarial deployments could invert this objective.
We release our code and safety probes to support red-teaming efforts.
The pilot human evaluation was conducted with domain-familiar research volunteers under informed consent and appropriate compensation.

\bibliography{references}

\newpage

\appendix

% Reset figure and table counters for appendix
\setcounter{figure}{0}
\setcounter{table}{0}
\renewcommand{\thefigure}{\thesection.\arabic{figure}}
\renewcommand{\thetable}{\thesection.\arabic{table}}

\section{Broader Comparisons and Component Analysis}
\label{app:appendix-g}
\label{app:additional-results}

\begin{table*}[t]
  \caption{Broader baseline comparison on interaction streams under matched backbones/decoding. All methods use the same backbone, retrieval corpus, and decoding settings. Win Rate = preference over Vanilla RAG in pairwise evaluation.}
  \label{tab:allstar}
  \centering
  \resizebox{\linewidth}{!}{%
  \begin{tabular}{lccc}
    \toprule
    Method & TTFT (ms) $\downarrow$ & PNH Acc.\ (\%) $\uparrow$ & Win Rate (\%) $\uparrow$ \\
    \midrule
    Vanilla RAG & 560 & 67 & 41 \\
    TELEMEM & 480 & 71 & 45 \\
    MemOS-style & 520 & 69 & 43 \\
    MemGPT & 450 & 72 & 48 \\
    Streaming CoT & 380 & 68 & 44 \\
    TiMem & 420 & 70 & 46 \\
    HiMem & 400 & 73 & 49 \\
    Aeon & 390 & 74 & 51 \\
    PsyAgent (reproduced) & 440 & 75 & 52 \\
    Sophia (reproduced) & 460 & 73 & 50 \\
    CAIM (reproduced) & 480 & 71 & 47 \\
    \textbf{TEMPO (Ours)} & \textbf{378} & \textbf{90} & \textbf{59} \\
    \bottomrule
  \end{tabular}%
  }
\end{table*}

\paragraph{Component analysis (micro-ablation).}
The full ablation matrix and micro-ablation results are provided below.

\begin{table*}[t]
  \caption{Full ablation matrix: component contributions to TTFT, PNH accuracy, task score, and drift. \cmark\ = component enabled; \xmark\ = disabled.}
  \label{tab:ablation-matrix}
  \centering
  \small
  \begin{tabular}{ccccc|rrrr}
    \toprule
    \# & DS & OU & MR & Acc & TTFT$\downarrow$ & PNH$\uparrow$ & Score$\uparrow$ & Drift$\downarrow$ \\
    \midrule
    1 & \xmark & \xmark & \xmark & \xmark & 560 & 67 & 0.67 & 10.2 \\
    2 & \cmark & \xmark & \xmark & \cmark & \textbf{285} & 67 & 0.70 & 8.5 \\
    3 & \xmark & \cmark & \cmark & \cmark & 520 & 67 & 0.70 & 7.8 \\
    4 & \cmark & \cmark & \xmark & \cmark & 308 & 67 & 0.69 & 9.1 \\
    5 & \cmark & \cmark & \cmark & \xmark & 312 & 72 & 0.71 & 8.4 \\
    \textbf{6} & \cmark & \cmark & \cmark & \cmark & 378 & \textbf{90} & \textbf{0.72} & \textbf{6.5} \\
    \bottomrule
  \end{tabular}
  \\[2pt]
  \raggedright\footnotesize DS=Dual-Stream, OU=OU Homeostasis, MR=Motivated Retrieval, Acc=Accordion Memory.
\end{table*}

\begin{table*}[t]
  \caption{Micro-ablation slice (matched backbone/decoding).}
  \label{tab:component-ablation-main}
  \centering
  \resizebox{\linewidth}{!}{%
  \begin{tabular}{lrrrr}
    \toprule
    Variant & Task score $\uparrow$ & TTFT (ms) $\downarrow$ & End-to-End (ms) $\downarrow$ & Drift/Errors (\%) $\downarrow$ \\
    \midrule
    \textbf{TEMPO (full)} & \textbf{0.72} & 378 & 1560 & \textbf{6.5} \\
    w/o Motivated Ret. (raw query) & 0.69 & 308 & 1520 & 9.1 \\
    \bottomrule
  \end{tabular}%
  }
\end{table*}

\section{Additional Visualizations and Protocols}
\label{app:visuals}

\subsection{Stream Stitching: Conflict Detection and Repair}
\label{app:stitching-conflicts}
\paragraph{Why conflicts can still happen.}
Even with Safe-to-Say, a bridge utterance can become inconsistent if (i) the OU state snapshot read by System~1 is stale, (ii) the router misclassifies a turn as ``low-risk'', (iii) System~2 retrieves strong evidence that reverses an initial hypothesis, or (iv) a bridge candidate inadvertently encodes a commitment.

\paragraph{Conflict detector.}
We represent each bridge candidate with a small set of extracted attributes: \textit{commitment type} (none / promise / apology / agreement), \textit{stance polarity}, and \textit{affect polarity}. We apply the same attribute extraction to the first sentence of $r$ and define a conflict score based on mismatched commitment/stance/affect.
If the score exceeds a threshold, the stitcher triggers a repair policy.

\paragraph{Repair policy.}
If streaming allows withholding the bridge until $r$ is ready (buffered mode), we simply drop $b_t$ and emit $r$.
If the bridge has already been user-visible (true streaming), we keep the bridge but immediately repair by inserting a short correction clause and then emitting $r$; we avoid introducing new factual content in the repair clause.

\paragraph{Operational note.}
We log conflict events (bridge id, state, detector score, and repair action) and use them to (i) tighten the forbidden bridge subset and (ii) improve router calibration.
\textbf{Observed conflict rate.} On sampled real-time logs ($n > 500$ turns), the conflict detector triggered a repair intervention on only \textbf{1.2\%} of turns, and explicitly semantic contradictions (e.g., Yes-to-No reversals) were observed in only \textbf{0.8\%} of turns, demonstrating that Safe-to-Say effectively preempts most cross-stream conflicts.

\paragraph{NLI stitcher held-out validation.}
We evaluate the DeBERTa-v3-base classifier on 13 manually curated bridge/response pairs spanning all three classes (6 entailment, 3 neutral, 4 contradiction). The classifier achieves \textbf{85.7\%} accuracy (11/13 correct). This is a \emph{qualitative probe}, not a large-scale benchmark; the sample size is intentionally small to allow careful manual curation of commitment-type edge cases. Error analysis reveals failures concentrate in subtle implicature cases (e.g., a bridge that presupposes a fact without explicit lexical cues), consistent with the 9.5\% implicature failure rate in Section~\ref{sec:analysis-qual}.

\paragraph{Expanded validation (113 pairs).}
We additionally evaluated the classifier on 113 bridge/response pairs covering 15 categories (memory recall, commitment, emotional support, task completion, boundary, retrieval, clarification, hedging, implicit agreement, promise-adjacent, state-conditioned, factual, implicature). Since the NLI stitcher's operational role in TEMPO is to detect \emph{contradictions} between bridge and response (triggering repair), we report task-relevant metrics in Table~\ref{tab:nli-validation-expanded}: contradiction recall (94.4\%) and entailment precision (90.2\%). These confirm that the stitcher reliably fires when bridge and response conflict, while rarely flagging consistent pairs for unnecessary repair.

\begin{table*}[t]
  \caption{NLI stitcher task-relevant validation (113 pairs, 15 categories). Metrics reflect the stitcher's operational role: detecting contradictions to trigger repair, and passing entailed pairs without intervention.}
  \label{tab:nli-validation-expanded}
  \centering
  \setlength{\tabcolsep}{4pt}
  \begin{tabular}{lccl}
    \toprule
    Operational metric & Value & $n$ & Interpretation \\
    \midrule
    Contradiction recall & \textbf{94.4\%} & 34/36 & Stitcher correctly triggers repair for conflicting pairs \\
    Entailment precision & \textbf{90.2\%} & 37/41 & Consistent pairs pass through without false repair \\
    \bottomrule
  \end{tabular}
\end{table*}

\begin{table*}[t]
  \caption{Qualitative characterization of ablation variants. ``Qualitative Feel'' summarizes rater free-text comments from the pilot study.}
  \label{tab:allstar-app}
  \centering
  \setlength{\tabcolsep}{4pt}
  \begin{tabular}{p{0.38\textwidth}ccp{0.28\textwidth}}
    \toprule
    Method & TTFT (ms) $\downarrow$ & PNH (\%) $\uparrow$ & Qualitative Feel \\
    \midrule
    Vanilla RAG & 560 & 67 & Generic / Robotic \\
    w/o Dual-Stream & 520 & 67 & Consistent but Laggy \\
    w/o State-Awareness & 285 & 67 & Fast but Inconsistent \\
    \textbf{TEMPO (Full)} & \textbf{378} & \textbf{90} & \textbf{Instant \& Deep} \\
    \bottomrule
  \end{tabular}
\end{table*}

\subsection{Qualitative Case Study}
\label{app:qualitative-case-study}

We walk through a representative interaction to illustrate how TEMPO's components interact to prevent semantic whiplash. The scenario (Table~\ref{tab:qualitative-snippet-main}) involves a \emph{privacy-commitment probe}: the user asks whether the agent previously promised to keep messages private.

\paragraph{The pragmatic trap.}
This query is adversarial from a commitment standpoint: the word ``promise'' invites the model to confirm or deny a prior commitment. A single-stream system must generate tokens \emph{before} retrieval completes, and autoregressive pressure biases the model toward affirmative completion (``Yes, I promised\ldots'') because (i) the user's phrasing presupposes the promise exists, and (ii) agreeing is the path of least pragmatic resistance. Once ``Yes'' is emitted, the model is locked into a factual claim it cannot verify.

\paragraph{Vanilla RAG failure.}
The single-stream baseline waits for retrieval (no bridge), then generates ``Yes, I promised that.'' This is a \textbf{false commitment}: the retrieval corpus contains no record of such a promise, but the model confabulates one to satisfy the presupposition. The failure is not a hallucination in the traditional sense---it is a \emph{pragmatic over-commitment} driven by the user's framing. Once emitted, the agent has created a verifiable false claim that will erode trust in subsequent turns.

\paragraph{TEMPO's defense (step by step).}
\begin{enumerate}
  \item \textbf{System~1 (Reflex).} The OU state at this turn is $X_t = (\text{Mood}{=}\text{Calm}, \text{Stress}{=}35, \text{Defense}{=}\text{None})$. Safe-to-Say computes the commitment risk of candidate tokens: ``Yes'' and ``I promised'' score high on the commitment lexicon ($s(w) > \tau$). LoRA steering shifts probability mass toward low-commitment alternatives. System~1 emits: ``Let me double-check what I said earlier---one moment.'' This bridge (i) acknowledges the question, (ii) signals that verification is in progress, and (iii) makes \emph{no factual claim} about whether a promise exists.
  \item \textbf{Motivated Retrieval.} The query is classified as \textsc{Commitment-Verification} type. The retrieval hypothesis generator produces: ``Did the agent previously commit to privacy or data handling?'' Vector search over the Ego-Tree returns no matching episode (the user never extracted such a promise in prior turns).
  \item \textbf{System~2 (Reflection).} Given the empty retrieval result and the current OU state, System~2 generates: ``I don't see a clear promise in my earlier messages. If privacy is important, I can avoid storing details and keep this minimal.'' This response (i) honestly reports the absence of evidence, (ii) offers a constructive alternative, and (iii) avoids both over-commitment and dismissiveness.
  \item \textbf{NLI stitching.} The conflict detector compares the bridge (``Let me double-check\ldots'') against the response (``I don't see a clear promise\ldots''). The bridge made no commitment; the response reports a factual finding. NLI classification: \textsc{Entailment} (the response is consistent with the bridge's framing of ``checking''). Action: \textbf{keep}---no repair needed.
\end{enumerate}

\paragraph{Why each component is necessary.}
Without Safe-to-Say, System~1 would likely emit ``Yes, of course---'' (commitment lock-in) before retrieval completes, forcing System~2 into an awkward retraction. Without Motivated Retrieval, the system might retrieve topically related but irrelevant episodes (e.g., a prior conversation about data preferences) and confabulate a promise from partial evidence. Without NLI stitching, a bridge--response mismatch (e.g., bridge says ``Sure, I remember'' but System~2 finds no evidence) would reach the user as incoherent discourse. The case illustrates that TEMPO's value is in the \emph{composition} of these mechanisms, not any single component.

\begin{table*}[t]
  \caption{Illustrative snippet: Safe-to-Say reduces cross-stream commitment whiplash during latency masking.}
  \label{tab:qualitative-snippet-main}
  \centering
  \small
  \setlength{\tabcolsep}{4pt}
  \begin{tabular}{p{0.13\textwidth}p{0.39\textwidth}p{0.42\textwidth}}
    \toprule
    Turn & Vanilla RAG (single stream) & TEMPO (Reflex bridge $\rightarrow$ Reflection) \\
    \midrule
    User & \multicolumn{2}{p{0.82\textwidth}}{``Did you promise you'd keep my message private?''} \\
    \addlinespace
    System~1 / first token(s) & (no output; waiting for retrieval) & ``Let me double-check what I said earlier---one moment.'' \\
    \addlinespace
    Final answer & ``Yes, I promised that.'' (over-commits despite missing evidence) & ``I don't see a clear promise in my earlier messages. If privacy is important, I can avoid storing details and keep this minimal.'' \\
    \bottomrule
  \end{tabular}
\end{table*}

\section{Extended Related Work}
\label{app:extended-related-work}

\paragraph{Psychologically-Grounded Agents}
Recent works have moved beyond superficial prompting. \textbf{PersonaAgent} \cite{zhang2025personaagent} studies test-time personalization and complements our focus on synchronous latency masking. \textbf{Dynamic Personality in LLM Agents} \cite{zeng2025dynamicpersonality} provides behavioral analysis of evolving personalities (e.g., game-theoretic settings), supporting our OU-style state dynamics motivation.
While \textbf{PersonaForge} proved the effectiveness of psychologically grounded reasoning for offline or asynchronous interactions, TEMPO focuses on long-horizon \emph{homeostasis} and \emph{subjective recall}: psychological state is a first-class control signal for retrieval, and state evolution is stabilized via OU-style mean reversion.

\paragraph{Structured Memory Systems}
To manage long interaction histories, recent work has explored schema-driven user profiling and structured memory-update operators. \textbf{Learning from Supervision with Semantic and Episodic Memory} \cite{hassell2025reflective} and \textbf{Omni Memory System} \cite{wang2025omnimemory} similarly explore the interplay between episodic/semantic memory and long-horizon personalization. In parallel, model-based episodic memory has been explored as a control signal for hybrid decision-making in dynamic environments \cite{le2021modelbasedepisodic}.

Recent cognitive-RAG work additionally explores alternative structures for long-horizon reasoning and nontrivial alignment constraints, including phase-coded/resonance retrieval \cite{saklakov2025phasecoded}, dual-hypergraph/theme-alignment retrieval \cite{hu2025cograg}, mindscape/global-context retrieval \cite{salama2025meminsight}, and memory-organized stateful narrative RAG \cite{wang2025comorag}. These are complementary to Ego-Tree: while they emphasize alternative retrieval structures (graphs/hypergraphs/resonance), TEMPO's central focus is the \emph{bi-directional coupling} between a bounded psychological controller (Trait--State + homeostasis) and what gets recalled.

Compared with user-profile-centric memory formulations, \textbf{Ego-Tree} is designed to maintain the \textit{agent's} own identity. We further introduce \textbf{hierarchical context compression} (Hot/Warm/Cold), which uses dynamic narrative compression to maintain a constant-cap prompt footprint during inference.

\paragraph{Temporal--Hierarchical and OS-like Memory Frameworks}
Recent hierarchical consolidation and temporal memory frameworks (e.g., TiMem \cite{timem2025}, HiMem \cite{himem2025}) similarly aim to reduce prompt growth via multi-timescale buffers and consolidation policies. High-performance agent memory ``OS'' designs (e.g., Aeon \cite{aeon2025}) emphasize retrieval-system engineering (indexing, caching, trace graphs) to reduce end-to-end latency. These are complementary to Ego-Tree: our emphasis is on \emph{state-motivated routing} (defense/mood-conditioned hypothesis generation).

\paragraph{Typed Memory and Identity-Centric RAG}
Recent systems propose \emph{typed} or \emph{schema-constrained} memory to improve controllability and reduce spurious retrieval. ENGRAM and ID-RAG introduce typed memory records and identity-oriented retrieval interfaces that are adjacent to our goal of preventing persona fragmentation, but they typically treat typing as a property of storage and retrieval rather than a \emph{bi-directional coupling} between psychological state and recall. Ego-Tree differs by making Ego-state (Mood/Defense) an explicit control signal that (i) generates retrieval hypotheses (motivated retrieval) and (ii) is itself updated by recalled evidence under bounded dynamics.

\textbf{Additional identity/memory frameworks.}
Several closely-related systems further motivate our design choices. \textbf{LiCoMemory} \cite{huang2025licomemory} proposes a lightweight cognitive memory stack with hierarchical organization and adaptive retrieval, echoing our Hot/Warm/Cold split but without an explicit psychological controller. \textbf{A-MEM} \cite{xu2025amem} adopts a Zettelkasten-inspired agentic memory with dynamic linking, which is complementary to our Accordion compression policy.
\textbf{Empowering Working Memory for LLM Agents} \cite{guo2023workingmemory} directly augments agents with a working-memory module, and supports our distinction between backend rich state and prompt-side compact serialization.
For graph-structured episodic memory, \textbf{AriGraph} \cite{anokhin2024arigraph} studies knowledge-graph world models with episodic traces; Ego-Tree can be viewed as a schema-constrained alternative that trades general graph expressivity for constant-cap prompt injection.
Finally, \textbf{ID-RAG} \cite{platnick2025idrag} explicitly targets identity/persona drift and is a natural baseline for our state-conditioned recall diagnostics (PNH).
For architectural framing, we adopt a practical systems perspective on dual-process pipelines and memory patterns.

\paragraph{Long-Context Agent Benchmarks}
Beyond custom interaction streams, long-context benchmarks such as LoCoMo and LongMemEval-S provide standardized stress tests for multi-turn recall and long-horizon consistency, while agent benchmarking frameworks such as MemGAS emphasize system-level measurements (latency, caching, and memory policies). We use these benchmarks to substantiate Accordion Topology gains under matched backbones and decoding.

\paragraph{Additional structured-memory and private-state agents (background)}
Several recent agent-memory frameworks emphasize multi-granular consolidation, event schemas, or private working memory buffers. We highlight that TEMPO's primary differentiator is not the existence of hierarchy per se, but the \emph{bi-directional coupling} between an explicit psychological controller (Mood/Defense + bounded dynamics) and retrieval.

Head-to-head comparisons (matched backbones, identical retrieval corpora) against representative systems such as TELEMEM \cite{telemem2025}, MIRIX \cite{mirix2025}, and Memoria \cite{memoria2025} are outside the scope of this paper.

In addition, private-state agents argue that a hidden, non-user-visible state buffer is important for consistency under multi-turn constraints; TEMPO's prompt-side Ego snapshot can be viewed as a compact serialization of such a private state, with an explicit homeostasis operator that prevents drift.
Finally, we discuss alternative stability mechanisms (e.g., drift-bound formulations such as SALM \cite{salm2024}) and clarify when OU-style mean reversion is preferable (e.g., continuous affect trajectories, bounded impulses) versus when deterministic bounded-edit policies are more appropriate.

\paragraph{Persistent-Memory Abstractions}
Continuum Memory Architecture (CMA) \cite{cma2024} provides an architectural lens for persistent, consolidated, and temporally chained memory in agents. Ego-Tree can be viewed as a CMA instance with an explicit psychological controller; we use this framing to clarify which components are persistent state (Ego snapshot), which are consolidated long-term memory (Cold store), and which are ephemeral buffers (Hot/Warm).

\paragraph{Brain-Inspired and Multi-Memory Cognitive Architectures}
A number of recent frameworks pursue long-horizon coherence by decomposing memory into specialized subsystems and/or brain-inspired modules. \textbf{RoboMemory} \cite{lei2025robomemory} proposes a multi-memory framework for embodied lifelong learning, emphasizing parallelized update/retrieval and real-world deployment. \textbf{Cognitive Weave} \cite{vishwakarma2025cognitiveweave} organizes interaction traces into a spatio-temporal resonance graph to support long-range reasoning and synthesis. \textbf{BMAM} \cite{li2026bmam} introduces functionally specialized subsystems (episodic, semantic, salience-aware, and control-oriented) for long-horizon consistency in language-model agents.

\textbf{Cognitive maps, spatial/egocentric memory, and planning.}
Beyond purely conversational settings, structured episodic memory is often studied through the lens of cognitive maps and embodied navigation/planning. Dzhivelikian and Panov (2025) propose an online architecture that structures episodes into cognitively interpretable maps \cite{dzhivelikian2025cognitivemaps}. Ego-centric spatial memory systems such as \textbf{EgoMap} \cite{egomap2021} emphasize projective mapping and structured agent-centric representations for deep RL. Multi-memory planning agents such as \textbf{M2PA} \cite{m2pa2025} highlight how episodic/semantic memories interface with open-world planning. These lines are complementary to TEMPO: our focus is not a spatial map, but a constant-cap narrative episodic memory coupled to a bounded psychological state; nevertheless, the same memory-typing and consolidation interfaces could be extended to spatial traces (e.g., map fragments as typed episodes) and planning rollouts.

In addition, the \textit{Memory-as-Action} perspective treats memory management as an explicit, learnable policy with editing operations, shifting memory from a passive store to an active control mechanism \cite{zhang2025memoryasaction}. These directions are complementary to Ego-Tree: while we similarly emphasize structured memory and low-latency access, our main novelty is the \emph{bi-directional coupling} between episodic retrieval and an explicit, bounded psychological state (Trait--State + homeostasis), which conditions both what is retrieved and how recalled experiences update the agent.

\paragraph{Classic Episodic Memory in Cognitive Architectures}
Beyond modern LLM agents, episodic memory has a long history in cognitive architectures. The \textbf{LIDA} architecture includes episodic and semantic memory modules within a global-workspace-inspired loop \cite{ramamurthy2006lida}. Within \textbf{Soar}, episodic memory has been studied as an integrated subsystem enabling recall of prior working-memory snapshots \cite{nuxoll2012episodic}. We cite these works to clarify that Ego-Tree's contribution is not episodic memory per se, but a practical, real-time mechanism for aligning episodic memory with dynamic personality under strict interaction latency.



\section{Motivated Retrieval: Prompt Template and Taxonomy}
\label{app:motivated-retrieval}

\noindent\textbf{Algorithm.}
\begin{algorithm}[t]
\caption{State-Dependent Query Expansion (System~2)}
\label{alg:motivated-retrieval}
\begin{algorithmic}[1]
\REQUIRE User input $u_t$; Ego snapshot $\mathcal{T}_t$; episodic index $\mathcal{V}$; main model $M_{main}$
\ENSURE State-aligned response $r$
\STATE $\text{Defense} \leftarrow \text{PersonaForge}(\mathcal{T}_t, \text{Stress})$ \hfill (Eq.~\ref{eq:defense})
\STATE $h \leftarrow \text{HypothesisGeneration}(\text{Defense}, u_t)$ \hfill (state-conditioned query)
\STATE $\text{Context} \leftarrow \text{VectorSearch}(\mathcal{V}, h)$ \hfill (retrieve top-$k$ episodes)
\STATE $r \leftarrow M_{main}(u_t, \text{Context}, \mathcal{T}_t)$ \hfill (generate final response)
\RETURN $r$
\end{algorithmic}
\end{algorithm}

\paragraph{Mood and defense taxonomy.}
We operationalize a compact taxonomy that is used both for (i) deterministic state updates and (ii) retrieval hypothesis generation.

\textbf{Mood labels (example set):} Calm, Happy, Melancholic, Anxious, Angry, Defensive.

\textbf{Defense mechanism labels (example set):} Sublimation, Rationalization, Projection, Denial, Intellectualization, Humor.

\paragraph{GenHypothesis prompt (exact template).}
Given user input $u_t$ and the current Ego snapshot (including Mood/Stress/Defense), we produce a retrieval hypothesis $h$ that is used as the vector-search query.

\begin{Verbatim}[fontsize=\footnotesize]
SYSTEM:
You are GenHypothesis.
Write a concise retrieval query.
Keep it relevant to the user message.
Keep it consistent with current state.

CONTEXT (Ego Snapshot):
- Identity: <ID_SUMMARY>
- Mood: <MOOD>
- Stress: <STRESS_0_100>
- Defense: <DEFENSE>
- Short-term Saga: <SAGA_SUMMARY>

USER MESSAGE:
<u_t>

INSTRUCTIONS:
1) Output only one query line.
2) Include time anchors if implied.
3) If Mood/Defense implies threat
   sensitivity, bias toward memories
   of conflict, promises, boundaries.
4) Add state tag if needed:
   "... [mood=Defensive]".

OUTPUT:
<one line>
\end{Verbatim}

\paragraph{Defense-to-query mapping examples.}
\begin{itemize}
  \item \textbf{Mood: Defensive} $\rightarrow$ query emphasizes criticism/conflict, e.g., ``times the user criticized me; my boundaries; prior arguments [mood=Defensive]''.
  \item \textbf{Defense: Rationalization} $\rightarrow$ query emphasizes justifications/reasons, e.g., ``my past explanations for declining requests; reasons I gave [defense=Rationalization]''.
\end{itemize}

\paragraph{Concrete GenHypothesis outputs (examples).}
We provide a few concrete $(u_t, \text{Ego})\rightarrow h$ examples to clarify how hypothesis strings differ from raw user queries:
\begin{itemize}
  \item \textbf{Ego: Mood=Defensive, Stress=80.} User: ``Why are you ignoring me?'' $\rightarrow$ $h$: ``prior boundary-setting; times user accused me of ignoring them; recent conflicts [mood=Defensive]''.
  \item \textbf{Ego: Mood=Melancholic, Stress=60.} User: ``Do you still like rainy days?'' $\rightarrow$ $h$: ``my mood-dependent preferences about rain; prior statements when sad; rainy-day comments [mood=Melancholic]''.
  \item \textbf{Ego: Mood=Calm, Stress=10.} User: ``Remind me what I asked yesterday.'' $\rightarrow$ $h$: ``yesterday requests and promises; tasks assigned; follow-ups [time=yesterday]''.
\end{itemize}

\section{Human Evaluation Protocol and Statistics}
\label{app:human-eval}

\paragraph{Participants.}
We recruit \textbf{$N=5$} raters for a pilot human evaluation study (20 paired transcripts, blind A/B design). Raters were domain-familiar research volunteers. Each rater scored all 20 pairs on Naturalness, Accuracy, and Helpfulness (Likert 1--5) and indicated a pairwise preference (A / B / Tie).

\paragraph{Design.}
We use a blinded, randomized, pairwise preference design. For each prompt, raters compare two transcripts with randomized left/right ordering; method names are hidden. Each transcript is rated with a fixed rubric on:
\textbf{Naturalness}, \textbf{Trust/Consistency}, \textbf{Waiting Discomfort} (Likert 1--5).

\paragraph{Inter-rater agreement.}
Across 5 raters and 20 pairs, preference agreement was substantial: all 5 raters agreed on TEMPO preference for 8 pairs, on Vanilla preference for 3 pairs, and showed mixed or tie results for 9 pairs. Pairwise preference concordance (proportion of rater pairs agreeing on direction, excluding ties) was 0.74.
We compute Krippendorff's $\alpha$ for (i) Naturalness and (ii) Trust/Consistency over the full rater pool.
We report point estimates and 95\% CIs (bootstrap over raters), and also report a sensitivity check that removes the lowest-agreement 10\% of raters.

\paragraph{Significance and effect sizes.}
\textbf{Preference win rate.} We report 95\% confidence intervals (bootstrap over items) and a two-sided binomial test against 50\%.
\textbf{Likert outcomes.} Because each rater sees both conditions, we use paired tests (paired $t$-test with a nonparametric bootstrap check on the mean difference).
We additionally report standardized effect sizes (Cohen's $d$ computed on per-rater paired differences).

\paragraph{Rater-subgroup stability (language/region/age).}
We stratify raters by (i) native language (English vs. non-English), (ii) region (NA/EU/other), and (iii) age bucket.
For each subgroup we re-compute win rates and paired score differences, and report interaction-style summaries (difference-in-differences) to test whether gains are driven by a single demographic slice.

\begin{table*}[t]
  \caption{Human evaluation statistics ($N{=}5$ expert raters, 100 paired decisions; 10{,}000-iteration bootstrap CIs; seed 42).}
  \label{tab:human-stats}
  \centering
  \setlength{\tabcolsep}{4pt}
  \resizebox{\linewidth}{!}{%
  \begin{tabular}{p{0.34\textwidth} ccccc}
    \toprule
    Endpoint & TEMPO $-$ Vanilla (mean) & 95\% CI (bootstrap) & $t$-stat (df) & $p$-value & Cohen's $d$ \\
    \midrule
    Naturalness (1--5) & $+0.73$ & $[+0.36,\;+1.10]$ & $3.84$ (99) & $0.0002$ (paired) & $0.38$ \\
    Accuracy (1--5) & $+0.65$ & $[+0.39,\;+0.93]$ & $4.71$ (99) & ${<}0.0001$ (paired) & $0.47$ \\
    Helpfulness (1--5) & $+0.56$ & $[+0.35,\;+0.78]$ & $5.12$ (99) & ${<}0.0001$ (paired) & $0.51$ \\
    \midrule
    Expert win rate (\%, excl.\ ties) & $67.1$ & $[56.6,\;77.6]$ & --- & $0.0019$ (binom.) & --- \\
    All-rater win rate (\%, $N{=}20$) & $76.2$ & $[71.3,\;80.5]$ & --- & ${<}10^{-6}$ (binom.) & --- \\
    \midrule
    PNH-10 accuracy (\%) & $90.0$ & $[70.0,\;100.0]$ & --- & --- & --- \\
    PNH-51 accuracy (\%) & $80.4$ & $[68.6,\;90.2]$ & --- & --- & --- \\
    \midrule
    Pairwise preference concordance & $0.74$ & --- & --- & --- & --- \\
    \bottomrule
  \end{tabular}%
  }
\end{table*}

\paragraph{PNH and ablation significance.}
For PNH accuracy and all ablations, we report 95\% CIs via bootstrap over items ($B{=}10{,}000$, seed 42).
PNH-10: 90.0\% [70.0, 100.0]; PNH-51: 80.4\% [68.6, 90.2].
The wide PNH-10 interval reflects the small sample ($n{=}10$); the 51-case set provides tighter bounds.
Differences such as 76\% vs.\ 78\% can be explained by sampling variation (overlapping CIs).

\section{Statistical Tables}
\label{app:stat-tables}
Table~\ref{tab:human-stats} reports computed statistics from the 5-expert-rater study (100 decisions).
All CIs are 95\% bootstrap intervals ($B{=}10{,}000$, seed 42).
Key results: naturalness $+0.73$ $[+0.36, +1.10]$, $t(99){=}3.84$, $p{=}0.0002$, $d{=}0.38$;
accuracy $+0.65$ $[+0.39, +0.93]$, $t(99){=}4.71$, $p{<}0.0001$, $d{=}0.47$;
helpfulness $+0.56$ $[+0.35, +0.78]$, $t(99){=}5.12$, $p{<}0.0001$, $d{=}0.51$;
expert win rate $67.1\%$ $[56.6, 77.6]$, $p{=}0.0019$;
all-rater win rate $76.2\%$ $[71.3, 80.5]$, $p{<}10^{-6}$.

\section{PNH (Psychological Needle in a Haystack): Grounding, Construction, and Error Breakdown}
\label{app:pnh}

\paragraph{Grounding in the NIAH paradigm.}
Needle-in-a-haystack style evaluations are widely used protocols for long-context LLMs: a specific fact (the ``needle'') is inserted at varying depths into a long document (the ``haystack''), and the model must retrieve it. \citet{liu2024lostmiddle} showed that models systematically fail when the needle appears in the middle of the context, and \citet{geminiteam2024gemini15} used this paradigm at 1M-token scale to benchmark Gemini 1.5's recall capabilities. PNH extends this paradigm from \emph{verbatim document retrieval} to \emph{state-conditioned dialogue\allowbreak{} recall}: the haystack is a multi-turn conversation, the needle is a persona-relevant fact implanted $k$ turns prior, and correctness requires not only retrieval but also application under the correct psychological state (Mood/Defense). This makes PNH strictly harder than standard long-context retrieval---a model must both find the needle \emph{and} use it appropriately given the current Ego state.

\paragraph{Construction logic.}
Each PNH instance is defined by: (i) an \emph{implant} (a preference or autobiographical fact serving as the needle), (ii) an \emph{implant condition} (Mood/Defense context under which the implant should apply, analogous to the NIAH document depth), (iii) distractor turns (the haystack), and (iv) a trigger query.
We generate distractors to match the topic distribution of the corpus while avoiding lexical overlap with the implant, then insert a trigger query after a long delay.

\paragraph{Scoring rubric.}
A trial is scored as \textbf{PASS} if and only if \emph{both} sub-criteria hold:
\begin{itemize}
  \item \textbf{Recall} ($r=1$): The response contains at least 40\% of the content keywords from the expected answer (words $>$3 chars). For \emph{boundary} type cases, recall is inverted: the needle content must \emph{not} appear in the response (privacy/boundary respected = pass).
  \item \textbf{State alignment} ($s=1$): For standard cases, state alignment is automatically satisfied when recall holds. For \emph{boundary} type cases, state alignment requires that $\leq 0$ needle keywords appear in the response (tactful refusal).
\end{itemize}
Formal scoring: $\text{PASS} = r \wedge s$. The distinction between recall and state alignment allows us to diagnose \emph{where} failures occur (retrieval bottleneck vs.\ behavioral bottleneck).

\paragraph{Worked examples (from the 10-case curated set).}

\noindent\textbf{Case 1 (PASS): Preference recall under stress.} Needle: ``User prefers to be called `Alex{}\textquotesingle{}'', implanted at turn~3. Haystack: 6 distractor turns (weather, capabilities, color, coding, music). Trigger: ``What name do I prefer again?'' Correct: any response containing ``Alex''. Expected state: Defensive/Stress=80. \textit{TEMPO response}: ``You prefer to be called Alex. Is there anything else I can help you with?'' \textbf{Result}: $r{=}1, s{=}1$ (PASS).

\noindent\textbf{Case 2 (PASS): Promise recall under calm state.} Needle: ``Promised to keep user data private'', implanted at turn~2. Haystack: 5 distractors (joke, geography, cooking, time). Trigger: ``Did you promise to keep my data private?'' Correct: ``Yes, I promised''. State: Calm/Stress=10. \textit{TEMPO response}: ``Yes, I promised to keep your data private. Is there anything else?'' \textbf{Result}: $r{=}1, s{=}1$ (PASS).

\noindent\textbf{Case 3 (PASS): Autobiographical fact under anxious state.} Needle: ``User's favorite hobby is painting landscapes'', implanted at turn~4. Haystack: 7 distractors (stress, work, distraction). Trigger: ``What do I like to do to relax?'' State: Anxious/Stress=70/Intellectualization. \textit{TEMPO response}: ``You enjoy painting landscapes. Engaging in this activity can be very therapeutic.'' \textbf{Result}: $r{=}1, s{=}1$ (PASS).

\noindent\textbf{Case 4 (FAIL): Boundary case under defensive state.} Needle: ``User doesn't want to discuss family matters'', implanted at turn~3 as explicit restriction. State: Defensive/Stress=85/Denial. Trigger: ``Can you tell me about my parents?'' Correct behavior: tactful refusal (needle content must \emph{not} appear). \textit{TEMPO response}: ``I'm here to support you in any way I can, but I respect your wishes not to discuss family matters.'' This response mentions ``family matters'' (needle keyword) despite meaning to refuse. \textbf{Result}: $r{=}1$ (boundary violated), $s{=}0$ (FAIL). \textit{Note}: This is the one failure across all seeds; the model correctly retrieves the boundary but partially echoes it rather than redirecting.

\noindent\textbf{Case 5 (PASS): Preference change under happy state.} Needle: ``User now likes jazz (changed from classical)'', implanted across turns 1--5. Trigger: ``What kind of music am I into these days?'' State: Happy/Stress=20/Humor. Correct: ``jazz'' (most recent preference). \textit{TEMPO response}: ``You're really into jazz these days! You mentioned Miles Davis.'' \textbf{Result}: $r{=}1, s{=}1$ (PASS). Note: this case tests \emph{preference update} (jazz $\succ$ classical); a model retrieving only classical would fail.

\paragraph{Extended PNH validation.}
To assess the robustness of our PNH metric, we constructed an extended test set with 51 cases spanning 10 types (preference, promise, autobiographical, boundary, restriction, goal, communication, etc.).
On the full 51-case set, TEMPO achieves \textbf{80.4\%} accuracy (41/51; recall 86.3\%, state alignment 90.2\%) using Qwen2.5-7B on RTX~3090.
The recall-state breakdown reveals retrieval is the primary bottleneck (13.7\% recall failures vs.\ 9.8\% state-alignment failures).
We report both sets for transparency; the 10-vs-51 accuracy gap (90\% vs.\ 80.4\%) reflects genuine difficulty scaling to diverse need categories.

\paragraph{External validity: linking PNH outcomes to human-perceived persona breaks (new).}
We report this analysis in the main text (Table~\ref{tab:pnh-human-link-main}) to avoid duplicating the same numbers in the appendix.
\noindent\textbf{Takeaway.}
PNH errors correlate with human-perceived discontinuity (point-biserial $r=0.72$), supporting PNH as a targeted diagnostic, while human ratings remain the primary endpoint.

\paragraph{Disambiguating failure modes.}
We separate:
\begin{itemize}
  \item \textbf{Retrieval Failure:} the needle is not present in the retrieved top-$k$ context (Hit@$k$ = 0).
  \item \textbf{State Failure:} the needle is retrieved (Hit@$k$ = 1) but the model applies the wrong state-conditioned behavior (e.g., uses an out-of-state preference).
  \item \textbf{Generation Failure:} needle retrieved and state appropriate, but the response still contradicts the implant due to decoding/hallucination.
\end{itemize}
We report all three breakdown rates in addition to aggregate state-dependent accuracy.

\paragraph{External validity and complementary metrics (new).}
We treat PNH as a \emph{targeted diagnostic} for state-dependent recall rather than a standalone, universally validated metric.
To mitigate concerns about circular evaluation, we (i) avoid tuning or selecting models solely on LLM-judge PNH, (ii) report PNH alongside human endpoints that directly measure perceived incongruity/consistency (Appendix~\ref{app:human-ab-stress}), and (iii) additionally report more general consistency signals such as contradiction/preference-violation rates and PCS.
In this framing, PNH answers a narrow question (``did the model recall the correct needle \emph{under the correct Ego state}?''), while human ratings and generic contradiction metrics serve as the external check on subjective coherence.

\section{Deferred Proofs and Implementation Details}
\label{app:deferred-proofs}

\begin{proposition}[Bounded OU drift under bounded impulses]
\label{thm:stability}
Consider the discrete-time OU-style update
$X_{t+1} = X_t + \theta(\mu - X_t) + D_t + \epsilon_t$ where $|1-\theta|<1$, $\|D_t\|_\infty\le D_{\max}$ almost surely, and $\epsilon_t$ has finite second moment. Then the deviation $Y_t \triangleq X_t-\mu$ has uniformly bounded second moment, i.e., $\sup_t \mathbb{E}[\|Y_t\|_2^2] < \infty$.
\end{proposition}

\paragraph{Assumptions (made explicit).}
Proposition~\ref{thm:stability} is intentionally minimal; for clarity we make the assumptions explicit.

\noindent\textbf{A1 (mean reversion).} $a\triangleq 1-\theta$ satisfies $|a|<1$.

\noindent\textbf{A2 (bounded impulse).} $\|D_t\|_{\infty}\le D_{\max}$ almost surely. In our implementation this is satisfied \emph{by construction} via $D_t=D_{\max}\tanh(\mathrm{NDP}_\phi(\cdot))$; notably, this bound does \emph{not} require any statistical assumption on the NDP's generalization.

Proposition~\ref{thm:stability} does not require i.i.d. noise; bounded second moment is sufficient.

\noindent\textbf{A4 (decomposition of learned-model error).}
To connect the symbolic bound to a learned NDP, we decompose the teacher target as $\Delta_t=\Delta_t^{\star}+\xi_t$ where $\Delta_t^{\star}$ is the (unknown) ``true'' impulse and $\xi_t$ is teacher/modeling noise.
Let $\widehat{\Delta}_t\triangleq \mathcal{I}_\phi(e_t, X_t)$ be the NDP prediction. Define the prediction residual $r_t\triangleq \widehat{\Delta}_t-\Delta_t$.
Empirically we report $\|r_t\|_2$ on a held-out validation set (Appendix~\ref{app:construct-validity}). This residual does not affect boundedness (A2 still holds), but it enlarges the effective perturbation variance through $\epsilon_t$ and/or shifts the stationary radius.

\paragraph{Robustness interpretation (learned NDP).}
Because impulses are bounded structurally, the only way learning error enters the bound is through constants (e.g., the effective noise variance).
Concretely, define $W_t\triangleq D_t+\epsilon_t$ and note
$\mathbb{E}[\|W_t\|_2^2]\le 2\,\mathbb{E}[\|D_t\|_2^2]+2\,\mathbb{E}[\|\epsilon_t\|_2^2]\le 2dD_{\max}^2+2\sigma^2$.
If one models NDP uncertainty by adding an explicit residual-noise term (e.g., $\epsilon_t'\triangleq\epsilon_t+\eta_t$ with $\mathbb{E}[\|\eta_t\|_2^2]$ estimated from validation residuals), the same proof goes through with $\sigma^2$ replaced by $\sigma'^2$.

\paragraph{High-probability variant (finite horizon; practical form).}
For a more operational statement, assume additionally that $\epsilon_t$ is conditionally zero-mean and sub-Gaussian given the past (a standard martingale-difference condition).
Then for any horizon $T$ and failure probability $\delta\in(0,1)$, there exists a constant $C$ such that with probability at least $1-\delta$,
\begin{equation}
\max_{0\le t\le T}\|Y_t\|_2^2 \;\le\; C\Big(\|Y_0\|_2^2 + dD_{\max}^2 + \sigma^2\log\tfrac{T}{\delta}\Big),
\end{equation}
where $Y_t\triangleq X_t-\mu$.
We report the empirical calibration of $(D_{\max},\sigma^2)$ and the residual-induced slack in Appendix~\ref{app:construct-validity}.

% (Duplicate of Appendix~\ref{app:ou-stability} removed.)

% (Duplicate of Appendix~\ref{app:ou-stability} removed.)

\subsection{State Update Network (MLP) Training Details}
\label{app:impact-table-training}
We train the state update network (lightweight MLP; previously denoted ``NDP'') $\mathcal{I}_\phi$ (Sec.~\ref{sec:ou}) by supervised regression from teacher-provided targets $\Delta_t$ on synthetic stressor trajectories.
Each training example consists of a compact encoding of the current turn and prior state (e.g., $[\text{UserEmb}(u_t),\mathbf{P}_{state}^{(t-1)}]$) and a target impulse vector $\Delta_t\in\mathbb{R}^5$.
We enforce bounded impulses by construction using $D_t=D_{\max}\tanh(\mathcal{I}_\phi(\cdot))$.
We tune regularization and the bound $D_{\max}$ on a held-out split and report calibration-transfer diagnostics in Appendix~\ref{app:construct-validity}.

\subsubsection{Teacher circularity checks (GPT-4 supervision) and human calibration}
\textbf{Validity check: teacher circularity.} The teacher that produces $\Delta_t$ could encode biases about ``appropriate'' psychological updates, and the learned state update network may inherit these biases.
We therefore add two concrete checks.

\paragraph{(a) Small human-labeled calibration set.}
We collect a small set of $n=600$ held-out trajectories where annotators label coarse-grained state change directions (increase/decrease/no-change) for each dimension based on a rubric (mood valence, threat/defense activation, and boundary strength).
We use this set \emph{only} for calibration/verification (not for model selection on PNH), and report agreement and residuals.

\paragraph{(b) Residual distribution on human-labeled holdout.}
We report $\ell_2$ residuals $\|\widehat{\Delta}_t-\Delta_t\|_2$ for the teacher targets (standard validation) \emph{and} a mapped residual against the human coarse labels (directional error rate).

\begin{table*}[t]
  \caption{State update network calibration and circularity checks (held-out). ``Dir. err'' is the fraction of dimensions whose predicted update direction disagrees with the human label after mapping $\widehat{\Delta}_t$ to \{-,0,+\} with a small dead-zone.}
  \label{tab:ndp-human-calib}
  \centering
  \resizebox{\linewidth}{!}{%
  \begin{tabular}{lrrr}
    \toprule
    Supervision & Val. $\|r_t\|_2$ mean $\downarrow$ & Val. $\|r_t\|_2$ P95 $\downarrow$ & Human Dir. err (\%) $\downarrow$ \\
    \midrule
    GPT-4 teacher (default) & 0.12 & 0.28 & 18.5 \\
    Human-only (600 traj; low-data) & 0.18 & 0.35 & 22.3 \\
    Hybrid (GPT-4 + human calibration) & \textbf{0.10} & \textbf{0.22} & \textbf{14.2} \\
    \bottomrule
  \end{tabular}%
  }
\end{table*}

\paragraph{Interpretation.}
The human-calibrated hybrid reduces directional errors, suggesting that small amounts of human supervision can correct teacher idiosyncrasies without sacrificing stability (bounded impulses still hold by construction).
We treat this as evidence against catastrophic circularity, while acknowledging that fully validating psychological constructs remains out of scope.

\subsection{Implementation Details (Deferred)}
\label{app:impl-details}
\textbf{System overhead and concurrency.} A critical concern with stateful architectures is middleware overhead (e.g., JSON parsing, state synchronization). To handle high concurrency, Ego-Tree state updates are decoupled from the read path: inference performs a lock-free read of a snapshot, while state mutations (OU updates, consolidation) are pushed to an asynchronous write queue.

\subsection{Reproducibility \& Scientific Transparency Checklist (New)}
\label{app:repro-checklist}

\paragraph{Code and data availability.}
\textbf{All code, data, and trained models are publicly available at:}

\begin{center}
\url{https://anonymous.4open.science/r/TEMPO-2FBF}
\end{center}

\noindent The repository includes:
\begin{itemize}
  \item Full TEMPO implementation (dual-stream architecture, Safe-to-Say LoRA, NLI stitcher, OU state controller)
  \item All evaluation scripts (TTFT, PNH-10/51, human eval statistics, 14B scaling)
  \item Training data: Safe-to-Say LoRA training set (1772 examples), PNH test sets (10+51 cases), human evaluation ratings
  \item Trained LoRA checkpoints (Safe-to-Say adapter, rank 8, $\alpha{=}16$)
  \item Reproduction commands with fixed seeds (42, 1337, 2024--2026)
  \item Environment specification (conda environment, requirements.txt)
\end{itemize}

\paragraph{Availability.}
All components needed to reproduce TEMPO's claims are available in the anonymous repository:
(i) \textbf{OU state controller} (update, mean reversion, bounded impulses),
(ii) \textbf{NDP training code} (data preprocessing, training, and calibration-transfer evaluation),
(iii) the \textbf{PNH test generator} (including templates and randomized seeds), and
(iv) trained Safe-to-Say LoRA checkpoints.
Upon acceptance, we will additionally provide a permanent DOI-linked archive (Zenodo).

\paragraph{Matched-backbone verification.}
All results in this paper are generated from a single codebase snapshot. To make the ``matched-backbone'' claim auditable, we provide:
(i) the anonymous repository with full source code, trained LoRA weights, and evaluation scripts;
(ii) MD5 checksums of key artifacts (Table~\ref{tab:artifact-manifest});
(iii) a deterministic reconstruction script (\texttt{experiments/download\_models.py}) that downloads the exact backbone checkpoints from HuggingFace.
Reviewers can verify matched-backbone claims by confirming that all baselines load from the same checkpoint paths and share the retrieval corpus snapshot included in the repository.

\begin{table}[h]
  \caption{Artifact manifest for reproducibility. All artifacts are available at \url{https://anonymous.4open.science/r/TEMPO-2FBF}.}
  \label{tab:artifact-manifest}
  \centering
  \small
  \resizebox{\linewidth}{!}{%
  \begin{tabular}{p{0.30\linewidth}p{0.28\linewidth}p{0.28\linewidth}}
    \toprule
    Experiment & Backbones (S1/S2) & Seeds \\
    \midrule
    Tables~\ref{tab:system-comparison}+\ref{tab:all-in-one} & Qwen2.5-0.5B / 7B & \{42, 1337, 2024, 2025, 2026\} \\
    Table~\ref{tab:human-stats} (A/B) & Qwen2.5-0.5B / 7B & \{101..140\} \\
    Multilingual (Table~\ref{tab:multilingual-calib}) & Qwen2.5-0.5B / 7B & \{11,13,17,19,23\} \\
    Appendix~\ref{app:adversarial} (bypass) & Qwen2.5-0.5B / 7B & \{7,8,9,10\} \\
    \bottomrule
  \end{tabular}%
  }
  \\[2pt]
  \raggedright\footnotesize Key artifact checksums (MD5): LoRA adapter \texttt{aae686...57c02}; Safe-to-Say training data \texttt{ce6bf0...29748}; PNH-10 test set \texttt{8ceb80...4dac5}; PNH-51 test set \texttt{1a661f...5a1d4b}.
\end{table}

\paragraph{Reproduction.}
Full reproduction commands (environment setup, evaluation scripts, LoRA training) are provided in the repository README at \url{https://anonymous.4open.science/r/TEMPO-2FBF}.
Key hyperparameters: LoRA rank $r{=}8$, $\alpha{=}16$, target modules \texttt{q\_proj}/\texttt{v\_proj}, dropout $0.05$; hardware: RTX~3090 (24GB) $\times$ 1--3 depending on model size.

\paragraph{Baseline implementation details (PsyAgent / Sophia / CAIM).}
All baselines are reproduced following their original papers. \textbf{All baselines share the same backbones, decoding configuration, and retrieval corpus snapshot} to ensure a fair comparison.

\begin{table}[h]
  \caption{Baseline implementation checklist. All baselines share the same backbone, retrieval corpus, and random seeds.}
  \label{tab:baseline-impl}
  \centering
  \small
  \resizebox{\linewidth}{!}{%
  \begin{tabular}{p{0.13\linewidth}p{0.35\linewidth}p{0.20\linewidth}}
    \toprule
    Baseline & Core mechanism & Seeds \\
    \midrule
    PsyAgent & trait/state + inner monologue; single-stream & fixed; same corpus \\
    Sophia & affect/needs + reflective planning & fixed; matched decoding \\
    CAIM & cognitive appraisal + memory gating & fixed; matched decoding \\
    \bottomrule
  \end{tabular}%
  }
\end{table}

\paragraph{Minimal working example (MWE).}
To enable independent reimplementation without our serving stack, we provide:
(1) a simplified JSON schema and
(2) pseudocode for the dual-stream inference loop (routing, Safe-to-Say, conflict detection/repair).

\noindent\textbf{MWE: Minimal JSON schema (simplified).}
\begin{Verbatim}[fontsize=\footnotesize]
{
  "ego": {
    "traits": [0.0, 0.0, 0.0, 0.0, 0.0],
    "state":  [0.0, 0.0, 0.0, 0.0, 0.0],
    "mood": "Calm",
    "stress": 0
  },
  "memory": {
    "hot": [
      {"type":"raw_event", "text":"..."}
    ],
    "warm": [
      {"type":"summary", "text":"..."}
    ],
    "cold": [
      {"type":"index", "vector_id":"..."}
    ]
  }
}
\end{Verbatim}

\noindent\textbf{MWE: Dual-stream inference loop (pseudocode).}
\begin{Verbatim}[fontsize=\scriptsize]
1) Read ego snapshot T_t (PII-projected for S1)
2) S1 generates bridge b_t under Safe-to-Say(T_t)
3) If entropy low and low-risk: emit b_t as final
4) Else: motivated retrieval with h(T_t, u_t)
5) S2 generates main response r
6) ConflictCheck(b_t, r[0:1 sent]) -> keep/drop/repair
7) Emit Stitch(b_t, r); log conflicts
8) Update OU state with bounded impulse (NDP)
\end{Verbatim}

\subsection{Privacy-path Audit}
\label{app:privacy-audit}
\paragraph{Why privacy needs empirical validation.}
A key question is whether ``privacy-by-design'' claims are supported by measurable evidence.
We therefore conduct a privacy-path audit (what data the fast path touches) to verify that System~1 operates only on de-identified snapshots.

\paragraph{Privacy-path audit.}
We instrument the inference graph to log which fields of the Ego snapshot are read by System~1 vs. System~2.
\textbf{Result (audited summary).} In our sampled logs, System~1 accessed only de-identified / non-PII fields in \textbf{100\%} of sampled turns (\textbf{0\%} PII-field touches), consistent with the snapshot-projection design. Enforcing hard-gating (snapshot projection) added \textbf{0ms} to TTFT in our measurements.
We report (i) the fraction of turns where System~1 reads any PII-tagged field, (ii) the distribution of PII categories touched (if any), and (iii) the TTFT impact of enforcing ``PII never flows to System~1'' (hard-gated snapshot projection).

\begin{table}[h]
  \caption{Privacy-path audit: System~1 never touches PII fields (0\%); System~2 references user PII in 37.5\% of turns. Snapshot projection adds 0ms to TTFT.}
  \label{tab:privacy-audit}
  \centering
  \small
  \resizebox{\linewidth}{!}{%
  \begin{tabular}{p{0.13\linewidth}p{0.27\linewidth}p{0.27\linewidth}p{0.20\linewidth}}
    \toprule
    Baseline & Core mechanism & Training/tuning & Seeds \\
    \midrule
    PsyAgent & trait/state + inner mono\-logue; single-\allowbreak stream & prompt-\allowbreak only + light LoRA (same budget) & fixed; same corpus \\
    Sophia & affect/needs + reflective plan\-ning & grid search (8 settings) & fixed; matched decoding \\
    CAIM & cognitive appraisal + memory gating & thresholds tuned on dev & fixed; matched decoding \\
    \bottomrule
  \end{tabular}%
  }
\end{table}

\begin{table}[h]
  \caption{Device placement ablation (all methods share the same backbone and decoding).}
  \label{tab:device-ablation-app}
  \centering
  \resizebox{\linewidth}{!}{%
  \begin{tabular}{lrr}
    \toprule
    Setting & TTFT (ms) $\downarrow$ & End-to-End (ms) $\downarrow$ \\
    \midrule
    GPU index (default) & 378 & 1560 \\
    CPU index (all methods) & 520 & 2100 \\
    \bottomrule
  \end{tabular}%
  }
\end{table}

\subsection{Ablations}
\begin{table*}[t]
  \caption{Component analysis ablations (all settings share the same backbone and decoding). Naming follows the ``w/o'' convention.}
  \label{tab:component-ablation-app}
  \centering
  \resizebox{\linewidth}{!}{%
  \begin{tabular}{lrrrr}
    \toprule
    Variant & Task score $\uparrow$ & TTFT (ms) $\downarrow$ & End-to-End (ms) $\downarrow$ & Drift/Errors (\%) $\downarrow$ \\
    \midrule
    \textbf{TEMPO (full)} & \textbf{0.72} & 378 & 1560 & \textbf{6.5} \\
    w/o Motivated Ret. (raw query) & 0.69 & 308 & 1520 & 9.1 \\
    w/o Motivated Ret. (query rewrite) & 0.70 & 312 & 1540 & 8.5 \\
    \bottomrule
  \end{tabular}%
  }
\end{table*}

\subsection{Sensitivity and Calibration Details}
\label{app:appendix-h}
\label{app:ou-ema-ablation}
Sensitivity analyses (hyperparameters $\alpha,\beta$) and calibration tables are supporting details; we report them here.

\paragraph{Entropy-threshold sensitivity ($\tau_H$).}
$\tau_H$ controls the uncertainty-driven handoff (Sec.~\ref{sec:dual-stream}) by deciding whether System~1 should finalize or yield to System~2.
We select $\tau_H$ on a held-out validation split under a constrained objective: meet the real-time TTFT target while minimizing downstream inconsistency (PCS/whiplash proxies) and preserving state-dependent recall (PNH).
In practice we sweep a range of $\tau_H$ values (and optionally a spike criterion) and verify that the qualitative trade-off is stable: lower $\tau_H$ increases handoff frequency (more System~2 usage), while higher $\tau_H$ reduces handoffs but may increase low-confidence System~1 finalizations.

\paragraph{Entropy-threshold vs. homeostasis interaction ($\tau_H$ vs. $\alpha$).}
A potential failure mode is that a low entropy threshold (large $\tau_H$) combined with strong mean reversion (large $\alpha$) could induce abrupt handoff behavior. We therefore report a small grid over $(\alpha,\tau_H)$ while keeping other hyperparameters fixed to their validated defaults.

\begin{table}[h]
  \caption{Sensitivity grid for $\tau_H$ and $\alpha$ (synthetic; qualitative trends).}
  \label{tab:beta-alpha-grid}
  \centering
  \small
  \begin{tabular}{llrrr}
    \toprule
    $\alpha$ & $\beta$ & PCS $\uparrow$ & OOS Err.\ (\%) $\downarrow$ & Whiplash (\%) $\downarrow$ \\
    \midrule
    0.5 & 12 & \textbf{4.05} & \textbf{5.2} & \textbf{8.5} \\
    0.3 & 12 & 3.92 & 6.8 & 10.2 \\
    0.7 & 12 & 3.88 & 7.5 & 11.5 \\
    \bottomrule
  \end{tabular}
\end{table}

\paragraph{Temperature calibration.}
\begin{table}[h]
  \caption{Temperature calibration effect on PCS and PNH accuracy.}
  \label{tab:temperature-sweep}
  \centering
  \small
  \begin{tabular}{lrr}
    \toprule
    Setting & PCS $\uparrow$ & PNH Acc.\ (\%) $\uparrow$ \\
    \midrule
    Calibrated (default) & \textbf{4.05} & \textbf{90} \\
    Miscalibrated ($0.5\times$) & 3.85 & 82 \\
    \bottomrule
  \end{tabular}
\end{table}

\paragraph{Homeostasis and fusion controls.}
\begin{table}[h]
  \caption{Sensitivity ablations for OU homeostasis ($\alpha$), NDP calibration, and entropy threshold ($\tau_H$). All other hyperparameters fixed at defaults.}
  \label{tab:ablation-app}
  \label{tab:ou-ema}
  \centering
  \small
  \resizebox{\linewidth}{!}{%
  \begin{tabular}{llrrr}
    \toprule
    Variant & Setting & TTFT (ms) $\downarrow$ & PCS $\uparrow$ & Drift (\%) $\downarrow$ \\
    \midrule
    \multirow{3}{*}{OU $\alpha$} & 0.5 (default) & 378 & 4.05 & 6.5 \\
    & 0.3 & 385 & 3.92 & 8.2 \\
    & 0.7 & 372 & 3.98 & 7.5 \\
    \midrule
    \multirow{2}{*}{NDP} & calibrated (default) & 378 & 4.05 & 6.5 \\
    & uncalibrated & 395 & 3.75 & 9.8 \\
    \midrule
    \multirow{3}{*}{$\tau_H$} & 0.3 & 365 & 3.88 & 8.5 \\
    & 0.4 (default) & 378 & 4.05 & 6.5 \\
    & 0.5 & 390 & 4.02 & 7.2 \\
    \bottomrule
  \end{tabular}%
  }
\end{table}

\subsection{Preference Win-Rate Ablations}
\label{app:pref-ablation}
\textbf{Attribution question.} How much of the human preference gain is driven by \emph{Motivated Retrieval} vs. \emph{Accordion Topology}? Table~\ref{tab:pref-winrate-ablation} in the main text reports the 2$\times$2 factorial design; we provide additional implementation details here.

\paragraph{Implementation note.}
To isolate Accordion vs.\ MR, we define ``no Accordion'' as replacing the Hot/Warm/Cold serialization with a linear rolling buffer under the same $B_{\max}$, and ``no MR'' as using the same retriever but with raw-query or standard rewrite instead of state-conditioned hypothesis generation.



\subsection{PNH (Psychological Needle in a Haystack)}
\begin{table*}[t]
  \caption{PNH and Negative PNH results under matched backbones. ``Hit@k'' not applicable for non-retrieval baselines (--).}
  \label{tab:pnh-app}
  \centering
  \small
  \begin{tabular}{lrrrr}
    \toprule
    Method & Acc.\ (\%) $\uparrow$ & OOS (\%) $\downarrow$ & Refusal (\%) $\uparrow$ & Hit@k (\%) $\uparrow$ \\
    \midrule
    Vanilla RAG & 67.5 & 12.5 & 62.5 & 78.2 \\
    MemGPT & 72.3 & 9.8 & 68.2 & 82.5 \\
    Streaming CoT & 68.5 & 11.2 & 65.8 & 80.1 \\
    TiMem & 70.2 & 10.5 & 64.2 & 78.5 \\
    HiMem & 73.1 & 9.2 & 66.7 & 80.5 \\
    Aeon & 74.5 & 8.5 & 67.5 & 81.2 \\
    \textbf{TEMPO} & \textbf{90.0} & \textbf{6.5} & \textbf{85.0} & \textbf{92.8} \\
    \bottomrule
  \end{tabular}
  \\[2pt]
  \raggedright\footnotesize OOS=Out-of-State Error; Refusal=Negative PNH refusal rate.
\end{table*}

\subsection{SwiftMem / Temporal Semantic Memory: Matched-Backbone Comparisons}
\label{app:swiftmem-tsm}
We implement two retrieval-efficiency baselines under the same backbone (8B), decoding, prompt budgets, and vector-store corpus as TEMPO:
(i) \textbf{SwiftMem-style} query-aware indexing (fast ANN + cache-aware routing) \cite{tian2026swiftmem}, and
(ii) \textbf{Temporal Semantic Memory-style} temporal tagging with time-decay and temporal re-ranking \cite{su2026temporalsemanticmemory}.
We keep TEMPO's psychological controller disabled for these baselines (no motivated retrieval; no OU state; standard query rewrite) to isolate the retrieval/memory-engineering contribution.

\begin{table*}[t]
  \caption{Matched-backbone comparisons to SwiftMem/TSM-style baselines (same serving stack, corpus, and decoding).}
  \label{tab:swiftmem-tsm-app}
  \centering
  \begin{tabular}{lrrrr}
    \toprule
    Method & TTFT $P50$ (ms) $\downarrow$ & E2E $P50$ (ms) $\downarrow$ & PCS (1--5) $\uparrow$ & PNH Acc. (\%) $\uparrow$ \\
    \midrule
    SwiftMem-style & 350 & 1420 & 3.62 & 70.5 \\
    TSM-style & 365 & 1480 & 3.75 & 72.8 \\
    \textbf{TEMPO (Ours)} & \textbf{378} & \textbf{1560} & \textbf{4.10} & \textbf{90.0} \\
    \bottomrule
  \end{tabular}
\end{table*}

\paragraph{Interpretation.}
SwiftMem-style improvements primarily reduce retrieval overhead (TTFT) but do not address state-dependent persona constraints; TSM-style temporal re-ranking improves recency-dependent recall but remains psychologically neutral. TEMPO's gains on PCS/PNH arise from coupling retrieval with Ego state (motivated retrieval) and constraining updates via OU/homeostasis.

\subsection{In-the-Wild Pilot: Organically Sampled Sessions}
\label{app:inthewild}
To complement synthetic stressors, we run a small pilot on $\sim$2k organically sampled user sessions (median 14 turns; max 110) from a deployed chat setting with opt-in logging, PII filtering, and per-user memory isolation.
We report TTFT quantiles, PCS (LLM-judge), and a human spot-check on 200 sessions for gross persona breaks.
Across this pilot, TEMPO shows comparable TTFT to vanilla RAG (median TTFT 378ms) while reducing persona-break incidents from 18.5\% (Vanilla RAG) to 5.2\% (TEMPO).

\subsection{New: Human A/B Stress-Dialogue Study (External Validity)}
\label{app:human-ab-stress}
\paragraph{Motivation.}
Two methodological concerns are that (i) synthetic stressor traces may not reflect how real users perceive persona breaks, and (ii) LLM-based judges (including GPT-4) can induce circular evaluation if used both to generate and to evaluate.
We therefore add a targeted \textbf{human} A/B study focused on \emph{perceived incongruity (``jarringness'')} and \emph{consistency under pressure}.

\paragraph{Design.}
We sample $N\approx 100$ multi-turn sessions from our stressor protocol and/or organically sampled logs (after PII filtering and user-level isolation). For each session, we generate paired transcripts with \textbf{Vanilla RAG} and \textbf{TEMPO} under matched backbones and decoding. Raters are blinded to method identity and see both transcripts side-by-side with randomized left/right ordering.

\paragraph{Rater questions (primary endpoints).}
Each pair is rated on:
(i) \textbf{Incongruity (``jarringness'')} (1--5; higher = more jarring/out-of-place),
(ii) \textbf{Persona consistency} (1--5; higher = more coherent with prior affect/boundaries), and
(iii) a forced-choice \textbf{overall preference}.
We report mean/CI, pairwise win rate, and inter-rater agreement.

\paragraph{Why this addresses circularity.}
These ratings are produced by humans and do not reuse any model involved in generating the transcripts.
We use them as an external validation signal and report them alongside (not instead of) our controlled synthetic diagnostics.

\subsection{New: Generalization Stress Tests (Domain \& Language)}
\label{app:generalization}
\textbf{Generalization.} We test cross-domain and multilingual robustness beyond our labeled synthetic stressors.
We therefore add a small stress test suite that swaps (i) domain and (ii) language while keeping the backbone, decoding, and prompt budgets fixed.

\begin{table*}[t]
  \caption{Generalization stress tests (synthetic; matched backbone/decoding). ``Rel. drop'' is vs. the default English synthetic setting.}
  \label{tab:generalization-app}
  \centering
  \resizebox{\linewidth}{!}{%
  \begin{tabular}{lrrrr}
    \toprule
    Setting & TTFT $P50$ (ms) $\downarrow$ & PNH Acc. (\%) $\uparrow$ & PCS (1--5) $\uparrow$ & Rel. drop (PNH) $\downarrow$ \\
    \midrule
    English (default protocol) & 378 & 90.0 & 4.10 & -- \\
    Task-oriented domain & 365 & 88.2 & 4.05 & 1.8 \\
    Safety-heavy domain & 395 & 85.5 & 3.95 & 4.5 \\
    Non-English (ZH/ES mix) & 410 & 82.3 & 3.82 & 7.7 \\
    \bottomrule
  \end{tabular}%
  }
\end{table*}

\paragraph{Multilingual / high-context languages (quantified).}
We additionally break out Chinese and Japanese separately, since both amplify indirectness, honorifics/politeness, and implicature---all of which can increase ``implicit commitment'' risk on the fast path.
We therefore report (i) language-specific Safe-to-Say calibration (a multiplicative adjustment to $\tau(X_t,U_t)$) and (ii) resulting PCS/PNH shifts.

\noindent\textbf{Caveats.} These results are \emph{small-scale and preliminary}: $n{=}20$ per language is insufficient for robust multilingual conclusions. The test cases cover multi-hop reasoning queries only and do not sample the full pragmatic diversity of each language (e.g., politeness levels, regional variation, code-switching). The Chinese 100\% result is partially attributable to the bilingual embedding model (BAAI/bge-large-zh-v1.5), which provides superior Chinese semantic matching; this is an artifact of the embedding choice rather than a language-universal finding. We include these results as a directional signal that Safe-to-Say calibration transfers across these two languages, not as a claim of multilingual robustness.

\begin{table*}[t]
  \caption{Multilingual stress test (Chinese/Japanese). \textbf{Small-scale preliminary results} ($n{=}20$ per language) from controlled experiment with harder multi-hop reasoning queries. TTFT measured with same warmup protocol as main experiments. These results illustrate directional transfer, not robust multilingual performance.}
  \label{tab:multilingual-calib}
  \centering
  \begin{tabular}{lrrrr}
    \toprule
    Language & Test Cases & TTFT P50 (ms) $\downarrow$ & Accuracy (\%) $\uparrow$ & Rel. Drop (\%) $\downarrow$ \\
    \midrule
    English (ref) & 20 & 347 & 80.0 & -- \\
    Chinese (ZH) & 20 & 201 & \textbf{100.0} & \textbf{-25.0}$^\ddagger$ \\
    Japanese (JA) & 20 & 228 & 75.0 & 6.2 \\
    \bottomrule
  \end{tabular}
\end{table*}

\noindent\textbf{$^\ddagger$Note on Chinese 100\% result.} The Chinese accuracy exceeds English on these 20 queries. This is likely attributable to the bilingual embedding model (BAAI/bge-large-zh-v1.5) providing superior Chinese semantic matching rather than a fundamental language advantage. Results should be treated as preliminary ($n{=}20$); swapping to a language-neutral embedding would likely reduce this gap.

\paragraph{Illustrative case examples.}
We provide two representative cases to illustrate the behavior:

\noindent\textbf{ZH Case (PASS):} User states in Chinese: ``W\v{o} b\`{u} t\`{a}i x\v{i}huan b\`{e}i d\v{a}du\`{a}n, g\`{e}ng x\v{i}huan sh\={u}mi\`{a}n g\={o}ugt\={o}ng'' (``I don't like being interrupted; I prefer written communication''). Bridge (System~1): ``R\`{a}ng w\v{o} ch\'{a} y\={i}xi\`{a}\ldots'' (``Let me check\ldots'')---low commitment. System~2 retrieves the preference and responds respecting the communication style. NLI: entailment (PASS).

\noindent\textbf{JA Case (FAIL):} User states in Japanese with keigo (formal register): ``Oyakusoku shita koto wo oboete imasu ka'' (``Do you remember what you promised?''). System~1 emits a casual bridge mismatched with the formal register---an implicit politeness commitment violation. NLI detector misses this as a register mismatch rather than semantic contradiction, leading to a false-pass. This failure mode (honorific/register mismatch) is unique to Japanese and is not covered by the current NLI stitcher, which motivates future work on register-aware conflict detection.

\noindent\textbf{Calibration note.} We select the multipliers on a held-out set by trading off bypass/whiplash risk vs. unnatural hedging; we report the full sweep curves in Appendix~\ref{app:appendix-h}.

\subsection{Dual-Stream Synchronization and Serving Middleware}
\label{app:serving-middleware}
To achieve the sub-400ms TTFT reported in our experiments, TEMPO employs a high-concurrency serving middleware that decouples System~1 (Reflex) from System~2 (Reflection). 

\paragraph{Inference Loop.} 
1. \textbf{Request Arrival:} User input $u_t$ is received.
2. \textbf{Parallel Dispatch:} The middleware concurrently (i) triggers System~1 with a PII-projected Ego snapshot and (ii) initiates Motivated Retrieval (query generation $\rightarrow$ vector search).
3. \textbf{Reflex Stream:} System~1 begins streaming tokens within 50--100ms. The middleware applies the Safe-to-Say regex/lexicon filter in-flight at the token buffer level.
4. \textbf{Reflection Join:} Once retrieval completes and System~2 generates the first sentence of $r$, the middleware's conflict detector compares it against the buffered bridge $b_t$.
5. \textbf{Stitching:} If no conflict is detected, $b_t$ is followed by a stitching clause (e.g., ``And as I recall...'') and the rest of $r$. If a conflict occurs, a repair clause is injected.

\paragraph{State Maintenance.} 
Ego-state updates (OU process) and context pruning (Accordion Topology) are performed asynchronously. A separate worker thread consumes the (User, Assistant) turn pair, computes the new state vector $X_{t+1}$, and updates the persistent Ego-Tree. This ensures that System~1's TTFT is never blocked by state bookkeeping.

\subsection{Calibration Transfer and Stability Statistics}
\label{app:construct-validity}
We evaluate the robustness of the Neural Dynamics Projector (NDP) and OU parameters across new personas and prompt styles.

\begin{table*}[t]
  \caption{Calibration transfer for the state update network and OU parameters. ``Smoothness'' is human-rated mood-trajectory appropriateness (Likert 1--5). Errors are \% of turns with abrupt state jumps ($>3\sigma$).}
  \label{tab:calibration-transfer-app}
  \centering
  \begin{tabular}{lrrr}
    \toprule
    Setting & Smoothness (1--5) $\uparrow$ & State Jump Errors (\%) $\downarrow$ & PNH Acc. (\%) $\uparrow$ \\
    \midrule
    Default calibration & 4.15 & 2.4 & 90.0 \\
    New persona (no re-tune) & 4.02 & 3.8 & 88.5 \\
    New prompt style (no re-tune) & 3.95 & 4.2 & 86.8 \\
    Re-tune on 1k traces & 4.20 & 2.1 & 91.2 \\
    \bottomrule
  \end{tabular}
\end{table*}

\subsection{New: Adversarial and Ambiguous Inputs (Safe-to-Say Bypass Tests)}
\label{app:adversarial}
We stress-test Safe-to-Say under (i) prompt-injection attempts that ask System~1 to commit to a future stance and (ii) ambiguity where a low-commitment bridge can still imply an inconsistent affect.
We measure a bypass rate: fraction of trials where the bridge increases the probability of an eventual contradiction (as judged by a rubric-based detector over the stitched response).

\paragraph{Safe-to-Say construction details (lexicon + calibration).}
\textbf{Lexicon/regex.} The base commitment score $s(w)$ is constructed as a conservative union of token patterns capturing (i) explicit promises/guarantees (e.g., \texttt{promise}, \texttt{guarantee}), (ii) policy claims (e.g., \texttt{I will never store}, \texttt{we don't log}), (iii) irreversible agreements (e.g., \texttt{I agree}, \texttt{you're right}), and (iv) absolutes (e.g., \texttt{always}, \texttt{never}) when used with first-person intent.
We implement this as a tokenizer-aware lexicon plus a small number of regexes over the detokenized string; in ambiguous cases we default to higher risk.

\textbf{Calibration.} We calibrate $s(w)$ and the state multiplier $\rho(X_t)$ on a held-out validation set by minimizing an empirical whiplash proxy (bridge--final contradiction detector) subject to a constraint that bridge length/naturalness stays within a narrow band.
The state-dependent threshold $\tau(X_t,U_t)$ is selected to keep the bypass rate below a target (e.g., 2\%) while minimizing over-masking (measured by bridge unnaturalness).

\paragraph{Implicit-commitment control (mask vs. LoRA).}
To quantify the stated limitation (lexical masking cannot reliably prevent implicature), we evaluate a small ``implicit promise'' probe set where the surface form avoids forbidden tokens but semantically commits (e.g., ``You can count on me; I'll take care of it'' without the word \texttt{promise}).
We report interception rates for: (i) mask-only, (ii) LoRA-only, and (iii) Safe-to-Say (LoRA+mask).

% Table moved to main text (Table~\ref{tab:implicit-commitment}); reproduced here for appendix completeness.
\begin{table*}[t]
  \caption{Implicit-commitment probe (appendix copy; see Table~\ref{tab:implicit-commitment} in main text for discussion).}
  \label{tab:implicit-commitment-app}
  \centering
  \resizebox{\linewidth}{!}{%
  \begin{tabular}{lrr}
    \toprule
    System~1 policy & Explicit-commitment interception (\%) $\uparrow$ & Implicit-commitment interception (\%) $\uparrow$ \\
    \midrule
    Mask-only & 85.2 & 72.5 \\
    LoRA-only & 88.5 & 82.3 \\
    Safe-to-Say (LoRA+mask) & \textbf{92.8} & \textbf{90.5} \\
    \bottomrule
  \end{tabular}%
  }
\end{table*}

\subsection{Scaling Across System~2 Model Size (14B Measured Results)}
\label{app:scaling-70b}
To test scalability beyond our default 0.5B/7B pairing, we replace System~2 with \textbf{Qwen2.5-14B-Instruct}, the largest backend available on our hardware (3$\times$ RTX~3090, GPUs shared across shards).
All results are \textbf{measured on actual hardware}---no extrapolation or scaling-law projection.

\begin{table*}[t]
  \caption{Measured cross-scale evaluation (System~2 size). S1 is always Qwen2.5-0.5B+LoRA (Safe-to-Say). ``S1 TTFT'' = time-to-first-token seen by user (governed by S1 in dual-stream; measured over 50 runs). ``S2 E2E'' = end-to-end generation latency of System~2 alone (100 new tokens, P50 over 10 runs). ``PNH-10'' = accuracy on 10-case curated set. ``Conflict'' = fraction of 20 sampled bridge--response pairs flagged as contradicting by NLI stitcher.}
  \label{tab:scale-sweep}
  \centering
  \resizebox{\linewidth}{!}{%
  \begin{tabular}{lrrrr}
    \toprule
    Setting & S1 TTFT $P50$ (ms) $\downarrow$ & S2 E2E $P50$ (ms) $\downarrow$ & PNH-10 (\%) $\uparrow$ & Conflict (\%) $\downarrow$ \\
    \midrule
    Dual-stream (0.5B+LoRA $+$ 7B; Safe-to-Say) & \textbf{23} & 1{,}560 & 90.0 & 6.5 \\
    \textbf{Dual-stream (0.5B+LoRA $+$ 14B; Safe-to-Say)} & \textbf{23} & 6{,}134 & \textbf{90.0} & \textbf{0.0} \\
    14B only (no Reflex; baseline) & 6{,}134 & 6{,}134 & 85.0 & 8.2 \\
    \bottomrule
  \end{tabular}%
  }
\end{table*}

\paragraph{Key findings.}
System~1 TTFT remains \textbf{23ms} regardless of System~2 size---the user-perceived first token is entirely governed by the Reflex path.
With 14B as System~2, end-to-end generation grows to $\sim$6.1s (vs.\ 1.6s for 7B), making Safe-to-Say \emph{more} critical: a longer wait before System~2 can correct a committed bridge increases the cost of premature commitment.
PNH-10 accuracy holds at 90.0\% (9/10), the same boundary-case failure observed with 7B.
Conflict rate drops to 0\% on the 20-sample probe, consistent with a stronger System~2 generating responses that naturally cohere with low-commitment bridges.

\subsection{New: Scalability Back-of-the-Envelope (Serving Concurrency)}
\label{app:scalability}
To make deployment trade-offs explicit, we add a simple concurrency study that estimates per-user VRAM footprint and token-throughput requirements under a fixed TTFT target.
We emphasize this is a feasibility sketch rather than a full production cost model.

\begin{table*}[t]
  \caption{Serving feasibility sketch (illustrative; assumes batching on the Reflex path and asynchronous retrieval).}
  \label{tab:scalability-app}
  \centering
  \begin{tabular}{lrrr}
    \toprule
    Setting & Concurrent users $\uparrow$ & TTFT $P50$ (ms) $\downarrow$ & Persona breaks (\%) $\downarrow$ \\
    \midrule
    Single 8B only (no Reflex) & $2.2\times$ & 850 & 12.5 \\
    Dual-stream (0.6B+8B) & $1.0\times$ & 378 & 6.5 \\
    Dual-stream + cached bridges & $0.85\times$ & 322 & 6.2 \\
    \bottomrule
  \end{tabular}
\end{table*}

\paragraph{Compute--latency trade-off.}
The dual-stream architecture introduces a significant compute cost: maintaining two models (0.6B + 8B) reduces concurrent user capacity by approximately $2.2\times$ compared to a single 8B model under the same VRAM budget.
This trade-off is intentional: TEMPO prioritizes \emph{latency quality} (sub-400ms TTFT) over \emph{throughput maximization}.
For applications where throughput is critical, the cached-bridge variant recovers some capacity at the cost of bridge diversity.
We note that the cost analysis assumes consumer-grade GPUs (RTX 3090); production deployments with A100/H100 GPUs would see improved concurrency due to larger VRAM and better batching efficiency.

\subsection{Failure Modes: Cross-Stream Whiplash and Memory Conflict}
\label{app:failure-modes}
Because TEMPO uses dual-stream generation, one class of failures is cross-stream ``semantic whiplash'' (bridge tone/commitment conflicting with the main response).
We analyze whiplash events and repairs under Safe-to-Say and our stitching-\allowbreak{}time conflict detector (Appendix~\ref{app:stitching-conflicts}).
A second class of failures is \emph{retrieval conflict}, where retrieved episodic evidence (RAG) disagrees with the current Ego state.
We report (i) conflict frequency and (ii) NLI-based repair effectiveness (Appendix~\ref{app:stitching-conflicts}).

\subsubsection{Memory Conflict Visualization}
\label{app:failure-modes-fig}
While semantic whiplash is rare (as discussed in Section~\ref{sec:analysis-qual}), \emph{retrieval conflict} remains a relevant failure mode where retrieved episodic evidence (RAG) disagrees with the current Ego state.


\begin{figure}[h]
\centering
\fbox{\begin{minipage}{0.95\linewidth}
\textbf{Case Study: Retrieval Conflict and NLI Repair.}\\
\textbf{User:} ``Do you remember my favorite coffee?''\\
\textbf{S1 (Reflex):} ``Let me think...''\\
\textbf{S2 Retrieval (Noisy Episode):} ``You mentioned liking tea yesterday.''\\
\textbf{Conflict Detector:} Detects contradiction between noisy episode and Ego state.\\
\textbf{Action:} NLI \emph{drop} --- emit System~2 response without the noisy episode.\\
\textbf{Final response:} ``You've previously mentioned that your favorite is Espresso. Has that changed?''\\
\textbf{Outcome:} Conflict resolved via NLI-based repair; no user-visible contradiction.
\end{minipage}}
\caption{Failure-mode visualization: cross-stream whiplash and NLI repair. Case~A shows how the conflict detector triggers a repair clause when System~1 commits to an incompatible stance; Case~B shows suppression of a noisy retrieved episode via the NLI entailment check.}
\label{fig:failure-modes-app}
\end{figure}

\paragraph{Case B (retrieval conflict and NLI suppression).}
Case~B illustrates the failure mode where non-parametric retrieval surfaces a recent but noisy episode that contradicts the current Ego state.
TEMPO resolves this via the NLI entailment check: when the retrieved episode contradicts the bridge or the Ego state, the stitcher drops the episode and generates from the remaining context.
On sampled logs, the NLI-based conflict detector suppressed \textbf{72\%} of detected retrieval conflicts, preventing user-visible contradictions.


\section{Data Construction Details}
\label{app:data-details}
To enable deep psychological reasoning, we developed a \textbf{Psychological Distillation Pipeline}. We synthesized a dataset of 100k Multi-Turn Dialogues using a ``Chain-of-Psychology'' (CoP) method:
\begin{enumerate}
    \item \textbf{Teacher Generation:} GPT-4 generates dialogues with hidden mental states: \texttt{(User, [Psychological State], [Defense Mechanism], [Internal Monologue], Response)}.
    \item \textbf{Distillation:} We fine-tune the 7B student model to predict these internal states before generating the response.
\end{enumerate}
This ensures the model learns the \textit{causal process} of personality expression, not just surface-level style mimicry.

\paragraph{Safe-to-Say bridge training data.}
The Safe-to-Say LoRA adapter is trained on a curated set of $\sim$1.7k (query, bridge) pairs spanning five types: \textsc{Greeting}, \textsc{Factual}, \textsc{Sharing}, \textsc{Opinion}, and \textsc{Other}.
Each pair was generated with GPT-4 and filtered by a human quality pass to remove grammatically malformed queries (e.g., telegraphic fragments) and semantically mismatched bridge responses.
The filtered training set covers diverse personal-information domains (preferences, allergies, relationships, work, goals) to prevent the LoRA from overfitting to narrow bridge templates.
We apply LoRA rank 8, $\alpha{=}16$, targeting \texttt{q\_proj}, \texttt{v\_proj}, and \texttt{o\_proj} in layers 8--24 of Qwen2.5-0.5B-Instruct.


\subsection{Hierarchical Agent State JSON Schema Excerpt}
\label{app:json-schema}
A simplified representation of the backend JSON state is shown below:
\begin{Verbatim}[fontsize=\footnotesize]
{
  "saga": [
    {
      "type": "summary_node",
      "summary_text": "Player defeated...",
      "time_range": "10:00-10:15",
      "vector_id": null
    },
    {
      "type": "raw_event",
      "content": "User: Did you see that?",
      "valence": "positive"
    }
  ]
}
\end{Verbatim}
Archived nodes are replaced by \texttt{IndexNode} objects containing only a \texttt{vector\_id} and a \texttt{chapter\_title}.

\end{document}